% ===================================================================
% FST-III: Biologische Stabilität und Nash-Frustration
% Zieljournal: Journal of Theoretical Biology / BioSystems
% Version: 3.0 (Quellencheck-Korrekturen, 2026-05-14)
% ===================================================================

\documentclass[12pt,a4paper]{article}

% Packages
\usepackage[utf8]{inputenc}
\usepackage[T1]{fontenc}
\usepackage{lmodern}
\input{glyphtounicode}
\pdfgentounicode=1
\usepackage[ngerman]{babel}
\usepackage{amsmath,amssymb,amsthm}
\usepackage{physics}
\usepackage{graphicx}
\usepackage{booktabs}
\usepackage{tabularx}
\usepackage{array}
\usepackage[margin=2.5cm]{geometry}
\usepackage{natbib}
\usepackage{enumitem}
\usepackage[protrusion=true,expansion=false]{microtype}
\usepackage{xurl}
\usepackage[unicode=true]{hyperref}
\usepackage{cleveref}

\graphicspath{{../results/}{../fst_nash/results/}}
\setlength{\emergencystretch}{3em}
\newcolumntype{L}[1]{>{\raggedright\arraybackslash}p{#1}}
\newcolumntype{Y}{>{\raggedright\arraybackslash}X}
\makeatletter
\renewcommand*{\HyperDestNameFilter}[1]{de.#1}
\makeatother

% Theorem environments
\newtheorem{proposition}{Proposition}
\newtheorem{conjecture}{Konjektur}
\newtheorem{hypothesis}{Hypothese}
\newtheorem{definition}{Definition}
\theoremstyle{remark}
\newtheorem{remark}{Bemerkung}

% Title
\title{Funktionale Stabilitätstheorie III: Biologische Stabilität und Nash-Frustration}

\author{Lukas Geiger\\
\textit{Unabhängiger Forscher, Bernau, Deutschland}\\
ORCID: \url{https://orcid.org/0009-0005-7296-1534}}

\date{Juni 2026}

\hypersetup{
  pdftitle={Funktionale Stabilitätstheorie III: Biologische Stabilität und Nash-Frustration},
  pdfauthor={Lukas Geiger},
  pdfsubject={Biologische Stabilität über Skalen in der Funktionalen Stabilitätstheorie},
  pdfkeywords={evolutionäre Spieltheorie, Proteinfaltung, Genregulation, dissipative Adaptation, Free Energy Principle, MEPP, Boolesche Netzwerke}
}

\begin{document}

\maketitle

\begin{abstract}
Dieser Beitrag entwickelt einen konjekturalen Rahmen für biologische Organisation von der molekularen bis zur Ökosystem-Skala, begründet in thermodynamischer Spieltheorie und dem Maximum Entropy Production Principle (MEPP). Er untersucht, ob biologische Stabilitätsprobleme als dissipative Strukturen modelliert werden können, deren Attraktoren Nash-artige Fixpunktbeschreibungen zulassen, sobald Spieler, Strategien und Payoffs auf der jeweiligen Skala spezifiziert sind. Der Rahmen bringt Jeremy Englands Theorie der dissipativen Adaptation, Karl Fristons Free Energy Principle, Stuart Kauffmans Boolesche Netzwerkdynamik und die klassische evolutionäre Spieltheorie in eine Arbeits-Hierarchie. Auf jeder Ebene---von der Proteinfaltung über Genregulation, zelluläre Koalitionen bis zur Ökosystemdynamik---bietet spieltheoretische Stabilität eine gemeinsame Sprache für biologische Organisation. Die Hauptübergänge der Evolution, identifiziert von Szathm\'ary und Maynard Smith, werden als mögliche Grobkörnungsschritte gelesen, die diese Struktur über Skalen hinweg erhalten können. Wir führen das Konzept der \emph{Nash-Frustration} als Komplement zur energetischen Frustration in Proteinstrukturen ein und prüfen es vorläufig anhand von NMR-Chemical-Shift-Perturbation-(CSP)-Daten an einem einzelnen Protein (Spearman $\rho_S = 0{,}44$, $p = 0{,}033$, $n = 24$). Das normalisierte Residuenranking ist im geprüften lokalen $\eta$-Fenster unabhängig von~$\eta$; absolute Instabilitätsskalen einschließlich $\rho(J)$ bleiben Kalibrierungsziele. Der Rahmen formuliert testbare Hypothesen für Proteinfaltungslandschaften, Genregulationsnetzwerk-Attraktoren und Ökosystem-Resilienz. Kosmologische Erweiterungen werden kurz erwähnt, sind aber explizit spekulativ und werden in einen Begleitbeitrag ausgelagert; sie sind nicht Teil der hier entwickelten biologischen Theorie.
\end{abstract}

\textbf{Schlüsselwörter:} evolutionäre Spieltheorie, ESS, Proteinfaltung, Genregulation, MEPP, dissipative Adaptation, Free Energy Principle, Boolesche Netzwerke, Hauptübergänge, Grobkörnung

\clearpage

% ===================================================================
\section{Einleitung}
\label{sec:introduction}
% ===================================================================

Die Entstehung und Stabilisierung biologischer Komplexität stellt die moderne Wissenschaft vor die Herausforderung, die fundamentalen Gesetze der Physik und die dynamischen Strategien des Lebens zu verbinden. Der Rahmen der Funktionalen Stabilitätstheorie (FST-III) behandelt biologische Systeme nicht bloß als Produkte zufälliger Mutationen, sondern als Kandidaten dissipativer Strukturen, die innerhalb thermodynamischer Randbedingungen operieren~\citep{EcosystemMEPP}.

Dieser Beitrag untersucht die vorgeschlagene Verbindung zwischen Spieltheorie und dem Maximum Entropy Production Principle (MEPP) und erforscht, ob biologische Stabilität durch die Perspektive spieltheoretischer Attraktoren innerhalb thermodynamischer Einschränkungen verstanden werden kann. Der Schwerpunkt ist primär biologisch (Abschnitte~\ref{sec:thermodynamics}--\ref{sec:ecosystem}); kosmologische Erweiterungen werden in Abschnitt~\ref{sec:cosmology} kurz erwähnt, sind aber nicht Teil des Kernarguments.

Die klassische evolutionäre Spieltheorie, bahnbrechend von Maynard Smith und Price~\citep{MaynardSmithPrice1973,MaynardSmith1982}, revolutionierte unser Verständnis biologischen Verhaltens. Doch die Standardbehandlung beschränkt sich weitgehend auf die Verhaltensökologie. Dieser Beitrag untersucht einen breiteren Vorschlag: spieltheoretische Beschreibungen können auf mehreren Ebenen biologischer Organisation formuliert werden, und diese Beschreibungen verbinden die Biologie mit der Physik durch eine analoge Stabilitätslogik, wie sie in FST-I~\citep{FSTI} und FST-II~\citep{FSTII} entwickelt wurde.

% ===================================================================
\section{Strukturelle Ausrichtung: FST-III als Pattern-A-Instanz}
\label{sec:structural-alignment}
% ===================================================================

Der hier entwickelte biologische Stabilitätsrahmen (ESS/FEP) wird als \textbf{Pattern-A-artige strukturelle Analogie} innerhalb des FST-Ökosystems vorgeschlagen (siehe~\cite{FST-Unified2026}), nicht als formale Ableitung aus dem mathematischen Pattern-A-Programm. Die jüngste \textbf{eichtheoretische Reformulierung der Riemannschen Vermutung} (siehe~\cite{GeigerRH2026c}) liefert ein motivierendes strukturelles Vorbild für diese Architektur.

Innerhalb dieses hierarchischen Stabilitätsrahmens werden biologische Systeme wie folgt verstanden:
\begin{enumerate}[label=\alph*)]
    \item \textbf{Analytische Rang-Endlichkeit (Null-Tail):} Biologische Komplexität kann gegen die kombinatorische Explosion des ``toten'' Zustands eingeschränkt sein, weil bei gegebenen Umwelt-Randbedingungen nur eine begrenzte Menge von Genotyp-Konfigurationen gleichzeitig ESS- und FEP-Stabilität erfüllen kann. Diese Endlichkeit ist konjektural, nicht bewiesen; sie ist motiviert durch die Beobachtung, dass die Schnittmenge von ESS-Bedingungen (Stabilität gegen Invasion) und thermodynamischer Lebensfähigkeit (positive Dissipation) genügend viele unabhängige Einschränkungen auferlegt, um die Lösungsmenge auf eine diskrete Kollektion zu reduzieren. Das unorganisierte ``Rauschen'' wird durch die entropische Einschränkung gefiltert und kann robuste Persistenz unterstützen.
    \item \textbf{Endliche Negativität als Eichsignal (Strukturelle Analogie):} Genetischer Konflikt, Krebs und Wettbewerb spielen eine Rolle analog zur endlichen Negativität im arithmetischen Kern. Sie werden hier nicht als zufällige Defekte behandelt, sondern als Signale unvollständiger organisatorischer Koordination, die durch neue Kooperationsebenen (Hauptübergänge) reduziert oder reorganisiert werden können. Diese Analogie ist strukturell, nicht formal: die mathematischen Mechanismen unterscheiden sich, aber das logische Muster (Destabilisierungssignale treiben Reorganisation) wird geteilt.
    \item \textbf{Eingeschränkte funktionale Positivität (Biologische Attraktoren):} Die stabilen Attraktoren eines Gennetzwerks oder Ökosystems können als ``eichstabile'' Zustände dargestellt werden, in denen das Freie-Energie-Funktional ein eingeschränktes Minimum erreicht. In Analogie zu Rangreduktionsmechanismen im RH-Programm wird phänotypische Stabilität hier als Ergebnis einer endlichen Menge epigenetischer Einschränkungen interpretiert, die thermodynamische Lebensfähigkeit stützen.
\end{enumerate}

Diese Ausrichtung rahmt FST-III als strukturelle Analogie---nicht als formale Ableitung---zu ``Funktionale Positivität unter Einschränkung''.

\begin{definition}[Resolvent-stabiler Fixpunkt]
\label{def:resolvent-stable}
Ein Fixpunkt eines dynamischen Systems heißt \emph{resolvent-stabil}, wenn die Störungsanalyse zweiter Ordnung (Hessian-Eigenwerte oder Resolvent-Entwicklungskoeffizienten) in allen Nicht-Eich-Richtungen strikt positive Werte liefert---d.\,h.\ alle Störungsmoden orthogonal zu symmetrieerzeugten Entartungen sind stabilisiert. In biologischen Systemen entsprechen ``Nicht-Eich-Richtungen'' Störungsmoden, die beobachtbare phänotypische Eigenschaften verändern (z.\,B.\ Proteinkonformation, Genexpression), während ``Eichrichtungen'' neutralen Freiheitsgraden entsprechen, die den Phänotyp invariant lassen (z.\,B.\ synonyme Codon-Substitutionen, neutrale Konformationsfluktuationen innerhalb eines Beckens). Dies ist eine strukturelle Analogie---motiviert durch, aber nicht formal abgeleitet von---der spektralen Lückeneigenschaft des FST-RH-Frameworks~\citep{GeigerRH2026b}; die Anwendung auf biologische Systeme erfordert domänenspezifische Verifikation, welche Richtungen wirklich ``Eich'' sind.
\end{definition}

% ===================================================================
\section{Thermodynamische Grundlagen: MEPP und dissipative Adaptation}
\label{sec:thermodynamics}
% ===================================================================

\subsection{Leben als Nicht-Gleichgewichts-Phänomen}

Die Einordnung biologischer Prozesse in der Thermodynamik beginnt mit der Erkenntnis, dass Leben ein Phänomen weit vom thermodynamischen Gleichgewicht ist~\citep{EcosystemMEPP}. Während geschlossene Systeme unvermeidlich zum Zustand maximaler Entropie und Wärmetod tendieren, erhalten sich biologische Systeme durch kontinuierlichen Export von Entropie an ihre Umgebung---ein Prozess, den Erwin Schrödinger 1944 als den Erwerb von ``Negentropie'' beschrieb~\citep{SchoolBiologicalPhysics,EcosystemMEPP,Kuzu2025}.

\subsection{Drei thermodynamische Hypothesen}

Innerhalb des FST-III-Rahmens beschreiben drei zentrale thermodynamische Hypothesen das Verhalten lebender Materie~\citep{EcosystemMEPP}:

\begin{enumerate}
    \item \textbf{Thermodynamischer Fingerabdruck des Lebens:} Lebende Gemeinschaften erhöhen die Rate der Entropieproduktion signifikant im Vergleich zu abiotischen Gegenstücken unter identischen Randbedingungen.

    \item \textbf{Zustandsselektionsprinzip (Stabilitätsinterpretation):} Wenn ein System mehrere stabile Zustände erreichen kann, tendiert die resolvent-stabile Konfiguration -- bei der Destabilisierungskanäle zweiter Ordnung unterdrückt sind -- dazu, hohe (wenn auch nicht notwendigerweise maximale) Entropieproduktionsraten aufzuweisen.

    \item \textbf{Gradientenantwortprinzip:} Wenn der äußere thermodynamische Gradient zunimmt, geht das System in einen neuen stabilen Zustand über, der durch noch höhere Entropieproduktion charakterisiert ist.
\end{enumerate}

\begin{table}[htbp]
\centering
\caption{Thermodynamische Prinzipien und ihre biologische Relevanz. MEPP gilt für stark getriebene Systeme; Prigogines minimale Entropieproduktion gilt nur nahe dem Gleichgewicht.}
\label{tab:thermo-principles}
\begin{tabularx}{\textwidth}{@{}L{0.22\textwidth} L{0.28\textwidth} Y@{}}
\toprule
\textbf{Prinzip} & \textbf{Beschreibung} & \textbf{Biologische Relevanz} \\
\midrule
MEPP & Stabilitätsfilter für hohe Dissipation & Metabolische Effizienzselektion \\
Zustandsselektion & Stabilität korreliert mit hohem $\dot{S}$ & Stabilisierung komplexer Nahrungsnetze \\
Gradientenantwort & Adaptation an zunehmenden Energieeintrag & Erhöhung photosynthetischer Effizienz \\
Prigogine (LMEP) & Minimale Entropie nahe Gleichgewicht & Nur für lineare, schwach getriebene Systeme \\
\bottomrule
\end{tabularx}
\end{table}

\begin{remark}[MEPP als Stabilitätsfilter, nicht als Selektionsgesetz]
\label{rem:mepp-stability-biology}
Gemäß der Resolventenanalyse aus dem Riemannschen-Vermutungs-Programm~\citep{GeigerRH2026b} sollten die drei obigen thermodynamischen Hypothesen als \emph{Stabilitätsfilter} gelesen werden, nicht als Selektionsgesetze. MEPP treibt biologische Systeme nicht kausal in Richtung Zustände maximaler Entropieproduktion. Stattdessen ist die Dynamik führender Ordnung entartet: viele biologische Konfigurationen sind in erster Ordnung lebensfähig. MEPP wirkt in dieser Lesart als Filter zweiter Ordnung, der resolventeninstabile Konfigurationen eliminiert---solche, die anfällig für störungsinduzierte Zusammenbrüche durch Kreuzkopplung zwischen metabolischen, strukturellen und replikativen Freiheitsgraden sind. Der Zustand mit ``höchster Entropieproduktion'' wird hier vorsichtiger als Kandidat eines resolventenstabilen Fixpunkts innerhalb der entarteten Mannigfaltigkeit thermodynamisch lebensfähiger Konfigurationen charakterisiert. Diese Verfeinerung adressiert die Gutachter-Kritik, dass MEPP als Selektionsgesetz verwendet wird: es ist ein Stabilitätskriterium, kein kausaler Treiber. Für die vollständige Entwicklung des MEPP-als-Stabilitätsfilter-Arguments im präbiotischen Kontext, einschließlich der formalen England-Ungleichung und der MEPP-Kontroverse, siehe FST-II~\citep{FSTII}; hier wenden wir dieses Prinzip spezifisch auf biologische Systeme an.
\end{remark}

\subsection{Dissipative Adaptation: Die Physik der Selbstorganisation}

Ein entscheidender Brückenschlag zwischen reiner Thermodynamik und biologischer Struktur wird durch Jeremy Englands Theorie der dissipativen Adaptation geliefert~\citep{England2013,EnglandNano}. Für die vollständige formale Herleitung von Englands Dissipationsschranke aus dem Crooks-Fluktuationstheorem, die spieltheoretische Lesung der Ungleichung und den Kausalrichtungsvorbehalt siehe FST-II~\citep{FSTII}; hier fassen wir nur die für den biologischen Kontext relevanten Aspekte zusammen.

England argumentiert, basierend auf Nicht-Gleichgewichtsstatistik, dass Materie, die einer starken externen Energiequelle ausgesetzt ist, dazu tendiert, sich in Konfigurationen zu organisieren, die besonders gut darin sind, Energie zu absorbieren und als Wärme zu dissipieren~\citep{England2013}.

Der Mechanismus der dissipativen Adaptation beruht auf der Irreversibilität von Zustandsübergängen. Wenn ein System Energie aus einer externen Quelle absorbiert, kann es energetische Aktivierungsbarrieren überwinden, die durch rein thermische Fluktuationen unüberwindbar wären. Sobald diese Energie dissipiert ist, fehlt dem System die Energie, um zurück zum Anfangszustand zu springen. Über die Zeit akkumulieren sich Strukturen, die hohe Resonanz mit der externen Energiequelle aufweisen~\citep{EnglandNano}.

Dieser Prozess erklärt die spontane Entstehung selbstreplizierender Strukturen: Replikation ist eine der effizientesten Methoden, um die Kapazität eines Systems zur Energieabsorption und -dissipation exponentiell zu steigern~\citep{EnglandNano}. In diesem Kontext erscheint ``Fitness'' nicht als rein biologisches Konzept, sondern als physikalische Eigenschaft struktureller Persistenz unter energetischem Antrieb~\citep{England2013}.

% ===================================================================
\section{Spieltheoretische Stabilität: Molekulare ESS}
\label{sec:molecular}
% ===================================================================

Sobald physikalische Strukturen die Schwelle zur biologischen Reproduktion überschreiten, werden ihre Interaktionen durch spieltheoretische Regeln bestimmt.

\subsection{Nash-Gleichgewicht und ESS}

Das zentrale Konzept ist das Nash-Gleichgewicht, 1950 von John Nash definiert, das Zustände beschreibt, in denen kein Akteur seinen Payoff durch einseitige Änderung seiner Strategie verbessern kann~\citep{NashPNAS,NashSFI}. In der Biologie entwickelten John Maynard Smith und George Price dies zur Evolutionarily Stable Strategy (ESS)~\citep{MaynardSmithPrice1973,MaynardSmith1982}.

\begin{definition}[Evolutionarily Stable Strategy]
Eine ESS ist eine Strategie, die, wenn sie von den meisten Mitgliedern einer Population angenommen wird, nicht von einer seltenen Mutanten-Strategie invadiert werden kann. Jede ESS ist ein Nash-Gleichgewicht, aber nicht jedes Nash-Gleichgewicht ist eine ESS, da ESS zusätzliche Stabilitätsbedingungen gegen Invasionen erfordert~\citep{MaynardSmith1982}.
\end{definition}

In FST-III wird dieses Prinzip auf die molekulare Ebene übertragen. Proteine, Enzyme und regulatorische RNAs werden als ``Spieler'' betrachtet, deren ``Strategien'' aus ihren strukturellen Konfigurationen und katalytischen Aktivitäten bestehen~\citep{GameBiochemistry}.

\begin{remark}[ESS als resolventenstabiler Fixpunkt]
\label{rem:ess-resolvent}
Eine mögliche Lesart aus der Resolventenanalyse der Riemannschen Vermutung~\citep{GeigerRH2026b} verfeinert die Interpretation von ESS. Eine ESS ist in dieser Lesart \emph{kein} globales Fitnessoptimum; sie ist ein Kandidat für einen resolventenstabilen Fixpunkt. In führender Ordnung ist die Fitnesslandschaft entartet---viele phänotypische Konfigurationen ergeben vergleichbare Payoffs. Die Selektion erfolgt in zweiter Ordnung durch resolventenartige Kopplungen zwischen den entarteten Modi: Invasionsdynamik, frequenzabhängige Interaktionen und räumliche Kopplung spalten die Entartung gemeinsam auf. Die ESS ist dann eine Kandidatenkonfiguration, bei der Destabilisierungskanäle zweiter Ordnung gleichzeitig unterdrückt sind. Nash-Frustration (Abschnitt~\ref{sec:molecular}) ist kein Defekt, sondern ein \emph{mögliches strukturelles Merkmal} dieser Resolventenarchitektur: sie reflektiert die residuale Kopplung zweiter Ordnung, die die Konfiguration stabilisieren und gleichzeitig eine Kristallisation in ein rigides, fragiles Optimum verhindern kann.
\end{remark}

\begin{table}[htbp]
\centering
\caption{Abbildung spieltheoretischer Konzepte auf physikalische und biologische Gegenstücke.}
\label{tab:game-biology-mapping}
\begin{tabularx}{\textwidth}{@{}L{0.18\textwidth} L{0.29\textwidth} Y@{}}
\toprule
\textbf{Spielkonzept} & \textbf{Physikalische Entsprechung} & \textbf{Biologisches Beispiel} \\
\midrule
Nash-Gleichgewicht & Stabiler stationärer Zustand & Phänotyp-Koexistenz \\
ESS & Attraktor mit lokaler Stabilität & Allel-Fixierung \\
Payoff-Matrix & Energie-/Fitnesslandschaft & Enzym-Substrat-Affinität \\
Gemischte Strategie & Populationsheterogenität & Arbeitsteilung in Kolonien \\
Nash-Frustration & Residuale Instabilität an der Stabilitätsgrenze & Funktionelle Flexibilität in Proteinen \\
\bottomrule
\end{tabularx}
\end{table}

\subsection{Proteinfaltung als Optimierungsspiel}

Ein markantes Beispiel ist die Proteinfaltung, die heuristisch als Spiel gegen die thermodynamische Landschaft modelliert werden kann~\citep{EnzymeFolding}. Die Aminosäuresequenz wird dabei nicht als intentionaler Akteur verstanden; die native Konformation wird im Modell als Nash-artiger Gleichgewichtskandidat in einem n-Personen-Spiel molekularer Kräfte repräsentiert:

\begin{itemize}
    \item \textbf{Spieler:} Einzelne Aminosäurereste
    \item \textbf{Strategien:} Rückgrat-Diederwinkel $(\phi, \psi)$
    \item \textbf{Payoff:} $-\Delta G$ (negative Freie-Energie-Änderung)
    \item \textbf{Nash-Gleichgewicht:} Native Konformation (Anfinsens Dogma~\citep{Anfinsen1973})
\end{itemize}

Fehlgefaltete Proteine (Prionen) repräsentieren alternative lokale Gleichgewichte---metastabile Konfigurationen, bei denen keine kleine einseitige Abweichung den Payoff verbessert, obwohl der globale Payoff suboptimal ist. Der Prionenfall ist besonders lehrreich: die fehlgefaltete Konformation ist kinetisch gefangen und verbreitet sich durch Template-gerichtete Konversion statt durch thermodynamische Selektion, was illustriert, dass Nash-artige Stabilität sowohl aus kinetischen Barrieren als auch aus energetischer Optimalität entstehen kann. Man beachte, dass die Prionenpropagation eine bekannte Ausnahme zu Anfinsens thermodynamischer Hypothese darstellt: der fehlgefaltete Zustand ist nicht das globale Freie-Energie-Minimum, sondern wird durch intermolekulares Templating und Aggregationskinetik stabilisiert.

\subsection{Nash-Frustration vs.\ energetische Frustration}

\noindent\textit{Notation:} $\rho(J)$ bezeichnet den Spektralradius der Jacobi-Matrix, $\rho_S$ den Spearman-Korrelationskoeffizienten und $\rho_i$ den normalisierten residuen-spezifischen Frustrationswert. Diese unterschiedlichen Verwendungen von $\rho$ folgen den jeweiligen Community-Konventionen und werden durchgehend durch Indizes disambiguiert.

Das Konzept der lokalen Frustration in Proteinstrukturen wurde extensiv durch den Protein Frustratometer~\citep{Ferreiro2007,Ferreiro2014} sowie proteom- und sequenzweite Deep-Learning-Modelle wie FrustraMPNN~\citep{Beining2026} und FrustrAI-Seq~\citep{Leusch2026} untersucht, der Reste identifiziert, bei denen die native Aminosäureidentität suboptimal ist, gegeben die lokale strukturelle Umgebung. Diese energetische Frustration korreliert mit funktionellen Stellen, Bindungsschnittstellen und allosterischen Regionen.

Der FST-III-Rahmen führt ein komplementäres Konzept ein: \textbf{Nash-Frustration}. Während energetische Frustration fragt ``welche Reste haben suboptimale lokale Energie?'', fragt Nash-Frustration ``welche Reste verhindern spieltheoretische Koordination?''

Formal wird \emph{residuen-spezifische Nash-Frustration} aus dem Best-Response-Jacobian $J = I - \eta H$ berechnet, wobei $H$ die $2n \times 2n$ Hesse-Matrix des Potentials ist (mit $n$ Resten, jeder mit zwei Diederwinkeln $\phi_i, \psi_i$) und $\eta > 0$ ein Lernraten-Parameter. Der Frustrationswert für Rest $i$ ist:
\begin{equation}
\text{frustration}(i) = \sum_{k: |\lambda_k| > 1} (|\lambda_k| - 1) \left( v_{k,\phi_i}^2 + v_{k,\psi_i}^2 \right)
\end{equation}
wobei $\lambda_k$ und $v_k$ Eigenwerte bzw.\ Eigenvektoren von $J$ sind. Die \emph{globale Nash-Frustration} des gesamten Proteins wird durch den Spektralradius $\rho(J) = \max_k |\lambda_k|$ zusammengefasst; wenn $\rho(J) > 1$, ist mindestens ein Modus unter Best-Response-Dynamik instabil. Diese beiden Maße---residuen-spezifische Scores und der globale Spektralradius---sind komplementär: $\rho(J)$ quantifiziert die Gesamtabweichung von Nash-Stabilität, während der residuen-spezifische Score die Instabilität lokalisiert.

\textbf{Koordinatenmetrik-Vorbehalt.}
Die obige Projektion verwendet die euklidische Norm der Eigenvektor-Komponenten im lokalen Diederwinkel-Chart $(\phi,\psi)$. Sie ist daher eine Chart-Diagnostik, kein intrinsisches koordinateninvariantes biomolekulares Observable. Eine intrinsische oder massengewichtete Variante sollte die quadrierte Komponentensumme durch eine metrikgewichtete Projektion mit der geeigneten Torsions-/Mannigfaltigkeitsmetrik ersetzen; der vorliegende Einzelprotein-Benchmark verwendet einen festen Chart und ist unter dieser Konvention zu lesen.

\textbf{Vorläufige Konsistenzprüfung und $\eta$-Unabhängigkeit.}
Wir prüften die Nash-Frustrations-Metrik gegen experimentelle NMR Chemical-Shift-Perturbation (CSP)-Daten, die aus BMRB-Eintrag~4428~\cite{BMRB4428} abgeleitet wurden. Dieser BMRB-Eintrag ist der HP67-Villin-Headpiece-NMR-Datensatz (primär PDB 1QQV); für den HP35-Strukturbenchmark (PDB 1YRF) wurde die dokumentierte Residuenabbildung vom C-terminalen HP67-Subdomänenbereich auf HP35 verwendet, wobei mutierte oder fehlende Positionen ausgeschlossen wurden. Die CSP pro Rest berechnet sich als $\text{CSP}_i = \sqrt{\Delta\delta_{\text{H}}^2 + (\Delta\delta_{\text{N}}/5)^2}$, wobei $\Delta\delta$ die sekundäre chemische Verschiebung relativ zu Random-Coil-Referenzwerten~\cite{Wishart1995} bezeichnet; der Faktor $1/5$ ist die Standard-Gewichtung für ${}^{15}$N~\cite{Williamson2013}.

Die Spearman-Rangkorrelation zwischen residuen-spezifischer Nash-Frustration und CSP beträgt $\rho_S = 0{,}44$ ($p = 0{,}033$, $n = 24$ Reste) und liefert ein erstes Konsistenzsignal dafür, dass Frustrationswerte die Rangordnung der NMR-beobachteten strukturellen Perturbation teilweise abbilden können. Die statistische Power dieses Tests ist moderat ($\approx 60\%$ bei $\alpha = 0{,}05$ für die beobachtete Effektstärke), und Replikation an weiteren Proteinen ist erforderlich. Die fünf am stärksten frustrierten Reste (M11, T12, S1, L0, L19) weisen CSP-Werte von 1{,}26, 0{,}68, 2{,}13, 1{,}44 und 1{,}10\,ppm auf. Ein Ausreißer (K23, CSP~$= 2{,}67$\,ppm bei niedriger Frustration) ist auf den bekannten Trp-Ringstromeffekt an dieser Position zurückzuführen.

Entscheidend ist, dass ein systematischer Scan über sieben $\eta$-Werte zwischen $0{,}001$ und $0{,}1$ zeigt, dass die Frustrations-CSP-Korrelation im geprüften lokalen Fenster \emph{unabhängig} von~$\eta$ ist (Variation~$< 10^{-6}$). Dies folgt aus der Struktur der Frustrationskarte, solange die instabilen Moden von negativen Hesse-Eigenwerten stammen ($h_k<0$) und keine positiven Überschussmoden ($h_k>2/\eta$) die Grenze $|\lambda_k|>1$ überschreiten: Dann skaliert $\eta$ das Instabilitätsgewicht linear, was sich unter Max-Normalisierung herauskürzt. Daher ist $\eta$ im geprüften Fenster kein Tuningparameter für das normalisierte Residuenranking; relevant bleibt es für absolute Score-Skalen, Spektralradius-Schwellen, Konvergenzdiagnostik und jedes Regime, in dem sich die aktive instabile Modenmenge ändert. Dieses Ergebnis, geprüft gegen zwei unabhängige Random-Coil-Referenzsätze (Wishart~1995, $r = 0{,}18$; Kjaergaard--Poulsen~2011~\cite{KjaergaardPoulsen2011}, $r = 0{,}20$; $|\Delta r| < 0{,}02$), stützt eine vorläufige Einzelprotein-Konsistenzdiagnostik, aber noch keinen extern kalibrierten Prädiktor.

\textbf{Absolutskalen-Gate (QUANT-III-2b).}
Ein erster Audit der absoluten Kalibrierung schließt diese Aufgabe nicht positiv. Die gespeicherten Score-Vektoren verwenden eine separate Maximalwert-Normalisierung pro Protein, $s_{p,i}=r_{p,i}/\max_j r_{p,j}$; alle 21 Protein--$\eta$-Zeilen des aktuellen Scans haben daher $\max_i s_{p,i}=1$. Da jede positive proteinspezifische Reskalierung $r_{p,i}\mapsto c_p r_{p,i}$ den Wert $s_{p,i}$ unverändert lässt, sind die rohe proteinübergreifende Skala und damit eine absolute Schwelle aus diesen Artefakten nicht identifizierbar. Bei der Software-Konvention $\eta=0{,}015$ erklärt eine gewöhnliche Kleinste-Quadrate-Abbildung der 24 HP35-Scores auf absolute CSP in-sample nur $R^2=0{,}034$. Die Leave-one-residue-out-Kalibrierung ergibt $Q^2=-0{,}176$ und einen mittleren absoluten Fehler von $0{,}496$\,ppm, schlechter als die entsprechende Trainingsmittelwert-Baseline ($0{,}474$\,ppm). Als separate Diagnostik mit funktionellen Residuenlabels statt CSP transferiert die vorhandene HP35-Schwelle auf Protein~G mit Balanced Accuracy $0{,}630$ und Sensitivität $0{,}417$; der umgekehrte Transfer ergibt $0{,}583$ bzw. $0{,}214$. Diese Prüfungen tragen keinen absoluten Kalibrierungsclaim. Die Reparatur verlangt unnormalisierte Residuenbeiträge samt Normalisierungsfaktoren, eine vorab festgelegte $\eta$-/Metrikkonvention und eine einmalige Kalibrierung an einem Protein gegen ein kontinuierliches Observable, gefolgt von der unveränderten Prüfung an einem unabhängigen Protein-Holdout (Skript: \texttt{quant\_iii\_2b\_absolute\_threshold\_gate.py}).

\begin{figure}[htbp]
\centering
\includegraphics[width=\textwidth]{eta_calibration/eta_calibration_csp.pdf}
\caption{Streudiagramm der residuenspezifischen Nash-Frustrationswerte gegen NMR-Chemical-Shift-Perturbation (CSP) für den HP35/1YRF-Benchmark, unter Verwendung von aus BMRB 4428 abgeleiteten HP67-Verschiebungen, die auf die C-terminale HP35-Subdomäne abgebildet wurden. Jeder Punkt repräsentiert einen Rest; die Spearman-Rangkorrelation beträgt $\rho_S = 0{,}44$ ($p = 0{,}033$). Die Farbkodierung zeigt den Lernraten-Parameter $\eta$ und demonstriert, dass die Frustrations-CSP-Beziehung im geprüften lokalen $\eta$-Scan invariant ist (alle $\eta$-Werte fallen auf dieselbe Kurve). Der Ausreißer K23 (hohe CSP, niedrige Frustration) ist auf den Trp-Ringstromeffekt zurückzuführen, nicht auf strukturelle Instabilität.}
\label{fig:eta-calibration}
\end{figure}

\begin{remark}[Berechnungsprotokoll]\label{rem:comp-protocol}
Die Nash-Frustrations-Analyse verwendet ein \emph{maßgeschneidertes spieltheoretisches Potential}, kein molekularmechanisches Kraftfeld (wie AMBER oder CHARMM).
Jeder Rest~$i$ wird als Spieler mit Strategie $\theta_i = (\phi_i, \psi_i)$ behandelt.
Das Potential $F = F_\phi + F_\psi$ besteht aus Einzelspieler-Termen---aminosäurespezifischen Ramachandran-Präferenzen $\mu_\phi^{(a_i)}, \mu_\psi^{(a_i)}$ mit Steifigkeitskonstanten~$k^{(a_i)}$---und paarweisen Kopplungstermen, gewichtet durch einen Radialbasisfunktions-Kern
\[
  w(r_{ij}) = \sum_{b=1}^{B} c_b \exp\!\bigl(-\tfrac{(r_{ij}-\mu_b)^2}{2\sigma^2}\bigr),
  \qquad B=12,\; \sigma = 1{,}0\,\text{\AA}.
\]
Kontakte werden durch einen C$_\alpha$--C$_\alpha$-Abstandscutoff von $r_\text{cut} = 8{,}0$\,\AA{} definiert.
Parameter (Ramachandran-Präferenzen, Steifigkeitskonstanten, RBF-Koeffizienten) werden aus der nativen PDB-Struktur durch regularisierte Kleinste-Quadrate-Anpassung rückentwickelt (Regularisierungsgewicht $\lambda_{\text{reg}} = 0{,}01$; Multi-Start-Gradientenabstieg mit 64 zufälligen Initialisierungen, je 6\,000 Schritte).

\textbf{Zirkularitätsvorbehalt.} Weil das Potential an die native Struktur angepasst wird, ist Konvergenz der Best-Response-Dynamik zu dieser Struktur erwartbar und stellt keine unabhängige Validierung dar. Der nichttriviale Gehalt der CSP-Korrelation ($\rho_S = 0{,}44$) liegt darin, dass die \emph{Rangordnung} der residuenspezifischen Frustration---abhängig von der Eigenvektorstruktur der angepassten Hesse-Matrix und nicht nur von der Fit-Güte---mit einem unabhängigen experimentellen Observable korreliert, das nicht in den Fit einging. Dennoch prägt das angepasste Potential die Hesse-Matrix und damit auch, welche Reste frustriert erscheinen. Ein stringenter Test des Nash-Frustrations-Konzepts erfordert daher entweder (a)~ein nicht an die native Struktur angepasstes physikalisches Potential, etwa ein All-Atom-Kraftfeld, oder (b)~Kreuzvalidierung, bei der an einem Protein angepasste Parameter Frustrationsmuster in strukturell unabhängigen Proteinen vorhersagen. Bis solche Tests vorliegen, ist das vorliegende Ergebnis als \emph{Konsistenzprüfung des selbstparametrisierten Modells} zu lesen, nicht als unabhängige Vorhersage.

Der Best-Response-Jacobian $J = I - \eta\, H$ verwendet eine Lernrate $\eta$; der Software-Standard ist $\eta = 0{,}015$, während die unten berichtete HP35-Pilotstudie $\eta = 0{,}01$ als absolute Skalenkonvention nahe $\rho(J) \approx 1$ verwendet. Die Hesse-Matrix~$H$ wird durch zentrale finite Differenzen berechnet ($\epsilon = 10^{-4}$).
Die Best-Response-Dynamik verwendet 300 Sweeps aus 16 zufälligen Startkonfigurationen.
Der Frustrationswert~$\rho_i$ projiziert die instabilen Eigenmoden von~$J$ (jene mit $|\lambda_k| > 1$) auf Rest~$i$ und normalisiert auf $[0,1]$.

Wie oben gezeigt, wird $\eta$ nicht optimiert, um die normalisierte CSP-Korrelation zu maximieren: Das rangbasierte Frustrationsmuster ist im geprüften Bereich $[0{,}001;\; 0{,}1]$ stabil, solange die aktive instabile Modenmenge gleich bleibt. Weitere Arbeiten sollten die absolute Score-Skala und $\rho(J)$-basierte Schwellen gegen unabhängige Observablen wie Wasserstoff-Deuterium-Austausch-Schutzfaktoren (McKnight et al.\ 1996~\cite{McKnight1996}) kalibrieren.
Quellcode: \texttt{protein\_fold\_nash\_pdb.py} (Python~3, BioPython, SciPy, NumPy).
Erweiterte Drei-Proteine-Analyse (LaTeX-Tabellenerzeugung, Konvergenzdiagnostik): \texttt{run\_extended\_analysis.py}.
\end{remark}

Tabelle~\ref{tab:nash-validation-ledger} trennt die empirischen und rechnerischen Evidenzblöcke von den jeweils nächsten Validierungsgates. Diese Trennung ist wichtig, weil der aktuelle Datensatz eine lokale, selbstparametrisierte Konsistenzprüfung stützt, während prädiktive Nutzung unabhängige Kontrollen und absolute Skalenkalibrierung verlangt.

\begin{table}[htbp]
\centering
\caption{Validierungs-Gate-Ledger für die Nash-Frustrations-Evidenz. Die Tabelle trennt aktuelle Observablen vom Mindestgate, das vor stärkeren prädiktiven Behauptungen nötig ist.}
\label{tab:nash-validation-ledger}
\footnotesize
\begin{tabularx}{\textwidth}{@{}L{0.21\textwidth} L{0.25\textwidth} L{0.22\textwidth} Y@{}}
\toprule
\textbf{Evidenzblock} & \textbf{Aktuelle Observable} & \textbf{Derzeitiger Status} & \textbf{Nächstes Gate} \\
\midrule
HP35/CSP-Rangprüfung & Aus BMRB 4428 abgeleitete CSP auf HP35 abgebildet; $\rho_S=0{,}44$, $p=0{,}033$, $n=24$ & Einzelprotein-Konsistenzsignal; kein extern kalibrierter Prädiktor & Replikation an weiteren Proteinen mit vorregistrierter Residuenabbildung und fester Skalenkonvention \\
$\eta$-Scan & Sieben Werte in $\{0{,}001,\ldots,0{,}1\}$; Variation des normalisierten Rankings $<10^{-6}$ & Lokale Ranginvarianz bei unveränderter instabiler aktiver Modenmenge & Absolute Scores und $\rho(J)$-Schwellen out-of-sample kalibrieren; aktive-Moden-Wechsel und positive Krümmungsüberschüsse verfolgen \\
Proteinübergreifende Konvergenz & HP35-gelernte Parameter auf 1PGA und 2XWR angewendet; alle Läufe konvergieren, aber $\phi_{\text{rms}}\approx1{,}25$--$1{,}30$ rad für Testproteine & Nur Konvergenzdiagnostik; kein Anspruch auf Strukturvorhersagegenauigkeit & 2XWR-Provenienz fixieren und reichere Potentiale mit Seitenketten-, Lösungsmittel- und Elektrostatiktermen prüfen \\
Benchmark funktioneller Reste & AUC 0{,}616 für HP35 und 0{,}688 für Protein~G & Schwaches bis moderates Signal über Zufall; unterhalb eigenständiger Diagnosestärke & Direkter Vergleich mit Ferreiro-Frustration, B-Faktoren/NMA, pLDDT und gematchten Sequenz-/Strukturkontrollen \\
Prospektive molekulare Tests & HDX-Schutzfaktoren, $\phi$-Werte, allosterische/Interface-Annotationen, Helix-Schmelzreihenfolge & Falsifizierbare Agenda spezifiziert; noch nicht ausgeführt & Multi-Protein-Tests mit vorregistrierten Schwellen durchführen und Negativkontrollen neben positiven Korrelationen berichten \\
\bottomrule
\end{tabularx}
\end{table}

\begin{table}[htbp]
\centering
\caption{Vergleich von energetischer Frustration (Ferreiro et al.) und Nash-Frustration (diese Arbeit). Die beiden Maße erfassen komplementäre Aspekte der Proteinstabilität.}
\label{tab:frustration-comparison}
\begin{tabularx}{\textwidth}{@{}L{0.18\textwidth} L{0.34\textwidth} Y@{}}
\toprule
\textbf{Konzept} & \textbf{Energetische Frustration} & \textbf{Nash-Frustration} \\
\midrule
Frage & Ist lokale Energie optimal? & Ist einseitige Abweichung stabil? \\
Methode & Kontaktenergie-Statistik & Best-Response-Jacobian-Eigenanalyse \\
Interpretation & Evolutionäre Einschränkung & Spieltheoretische Koordination \\
Hohe Werte zeigen & Funktionelle Stellen, Schnittstellen & Faltungsintermediate, kinetische Fallen \\
\bottomrule
\end{tabularx}
\end{table}

Rechnerische Experimente an drei strukturell verschiedenen Proteinen untersuchen, ob Nash-Frustration biologisch bedeutsame Reste lokalisiert und ob Best-Response-Dynamik jenseits der Trainingsstruktur konvergiert. Tabelle~\ref{tab:nash-results} fasst die Ergebnisse zusammen.

\begin{table}[htbp]
\centering
\caption{Nash-Gleichgewichtsanalyse von Proteinfaltungen. Parameter gelernt aus HP35 (Villin Headpiece, 1YRF) und angewendet, um Konvergenz der Best-Response-Dynamik für Protein G (1PGA, 56 Reste) und die hier genutzte p53-DNA-Bindungsdomänenstruktur (2XWR, 199 modellierte Reste) zu testen. Best-Response-Dynamik konvergiert über alle drei Proteine trotz mehrerer Proteingrößen; die großen $\phi_{\text{rms}}$-Werte zeigen jedoch, dass die konvergierten Strukturen den nativen Faltungen nicht eng entsprechen.}
\label{tab:nash-results}
\begin{tabular}{lccccc}
\toprule
\textbf{Protein} & \textbf{Reste} & $\phi_{\text{rms}}$ [rad] & $\psi_{\text{rms}}$ [rad] & \textbf{Konvergiert} & $F_{\text{best}}$ \\
\midrule
HP35 (1YRF, Training) & 35 & 1{,}054 & 1{,}661 & Ja (267 Sw) & $-4{,}61$ \\
Protein G (1PGA) & 56 & 1{,}300 & 2{,}002 & Ja (282 Sw) & $-10{,}61$ \\
p53 DBD (2XWR) & 199 & 1{,}247 & 2{,}114 & Ja & $-39{,}43$ \\
\bottomrule
\end{tabular}
\end{table}

\textbf{Referenzprotein (HP35).} Nash-Frustration konzentriert sich an Loop-Regionen (insbesondere Met11, Pro20 mit Werten von 1{,}0 und 0{,}49), während Helix-Kerne (Ser14, Asn18, Trp22) nahezu Null-Frustration aufweisen---konsistent mit ihrer Rolle als ``strukturelle Anker'' im Faltungsspiel. Der Spektralradius $\rho(J) = 1{,}013$ mit 80\% positiven Hessian-Eigenwerten und mittlerer Frustration 0{,}082 charakterisiert eine Nahe-Nash-Struktur. Dieses Muster unterscheidet sich von energetischer Frustration, die typischerweise an oberflächenexponierten funktionellen Resten ihren Höchstwert erreicht.

\textbf{Proteinübergreifender Konvergenztest.} Die aus HP35 gelernten Modellparameter wurden verwendet, um zu testen, ob Best-Response-Dynamik für Protein G (1PGA, 56 Reste) und die hier genutzte p53-Tumorsuppressor-DNA-Bindungsdomänenstruktur (2XWR, 199 modellierte Reste) konvergiert. Beide Läufe konvergierten zu wohldefinierten Fixpunkten (Tabelle~\ref{tab:nash-results}). Die p53-Domäne---ungefähr sechsmal größer als das Referenzprotein und einer anderen Faltungsklasse zugehörig (Immunglobulin-artiges $\beta$-Sandwich vs.\ $\alpha$-helikal)---erreichte einen $\phi_{\text{rms}}$ von 1{,}25\,rad, vergleichbar mit Protein G (1{,}30\,rad). Diese $\phi_{\text{rms}}$-Werte entsprechen mittleren Abweichungen von etwa 72$^\circ$ pro Rest und zeigen, dass die konvergierten Strukturen \emph{nicht} nahe an den nativen Faltungen liegen. Die Schlüsselbeobachtung ist daher auf \emph{Konvergenz} der Best-Response-Dynamik über Proteine verschiedener Größe und Faltungsklasse begrenzt, nicht auf strukturelle Vorhersagegenauigkeit. Die Potentialspiel-Formulierung erfasst Konvergenzeigenschaften der angepassten Dynamik; Strukturvorhersage würde ein reichhaltigeres Potential mit Seitenkettenwechselwirkungen, Lösungsmitteleffekten und Elektrostatik erfordern.

\paragraph{Rechnerische Beobachtung (Nash-Stabilitätsgrenze).}
Native Proteinfaltungen sind keine exakten Nash-Gleichgewichte unter lokaler Best-Response-Dynamik, sondern liegen nahe der Stabilitätsgrenze ($\rho(J) \approx 1$). Der Spektralradius $\rho(J) = 1{,}013$ für HP35 zeigt, dass die native Struktur nahezu selbststabilisierend ist, mit 14 von 70 Eigenwerten von $J$, die betragsmäßig die Eins überschreiten, und diese instabilen Moden sind nur an funktionellen Scharnierregionen lokalisiert, statt globalem Koordinationsversagen.

Diese Beobachtung legt nahe, dass Nash-Stabilität \textit{strikter} ist als Energieminimierung. Eine Struktur kann ein lokales Energieminimum sein, während sie dennoch einseitige Verbesserungen für einzelne Reste erlaubt. Die lokalisierte Nash-Frustration bei Met11 ist konsistent mit einem funktionellen Freiheitsgrad---einem ``Scharnier'', das konformationelle Flexibilität ermöglicht, während die Helix-Kerne stabile Koalitionen bilden. In der vorliegenden Interpretation balancieren funktionelle Proteine Stabilität gegen Flexibilität, anstatt in rigide Nash-Optima zu kristallisieren. Eine systematische Kalibrierung der absoluten Score-Skala gegen experimentelle Daten---etwa durch Abgleich der Nash-instabilen Reste mit bekannten allosterischen Stellen oder Faltungskeimen aus H/D-Austausch-Experimenten---wäre nötig, bevor dieser Machbarkeitsnachweis zu einem prädiktiven Werkzeug wird. Bis eine solche Kalibrierung durchgeführt ist, sollten die hier berichteten numerischen Werte ($\rho(J) = 1{,}013$, Frustrations-Maxima bei Met11 und Pro20) als parameterabhängige Illustrationen statt als robuste quantitative Vorhersagen interpretiert werden.
Im Sinne der FEP-MEPP-Kompatibilität (Konjektur~\ref{conj:fep-mepp}) kann Nash-Frustration $\rho(J)$ als Modelldiagnostik für lokale Spannung zwischen interner Freie-Energie-Minimierung (mEPP-Regime) und externer Entropieproduktions-Maximierung (MEPP-Regime) gelesen werden: Hohe Frustration markiert Reste, an denen der angepasste gefaltete Zustand dissipative Effizienz gegen strukturelle Stabilität abwägen könnte. Dies ist eine Stabilitätsinterpretation des Modells, keine direkte organismische Wahlhandlung.

\paragraph{Geplanter quantitativer Vergleich.}
Ein systematischer Vergleich der Nash-Frustrationswerte auf Residuen-Ebene mit den energetischen Frustrationsindizes des Protein Frustratometer~2.0~\cite{ParraFerreiro2016} stellt einen zentralen zukünftigen Validierungsschritt dar. Die Kernfrage ist, ob Nash-frustrierte Reste (hohe $|\lambda_k|-1$-Beiträge) mit energetisch frustrierten Resten überlappen, diese komplementär ergänzen oder orthogonal zu ihnen sind. Vorläufige Erwartungen legen partielle Komplementarität nahe: Energetische Frustration identifiziert evolutionäre Einschränkungen und Bindungsstellen, während Nash-Frustration kinetische Fallen und Faltungsintermediate hervorheben sollte (Tabelle~\ref{tab:frustration-comparison}). Ein Datensatz von Proteinen mit bekannten allosterischen Stellen und experimentell bestimmten $\phi$-Werten würde die empirische Grundlage für diesen Vergleich liefern. Vergleiche rangbasierter Karten erfordern kein optimiertes $\eta$; absolute Schwellen und globale $\rho(J)$-Vergleiche erfordern eine explizite Skalenkonvention oder eine Kalibrierung an unabhängigen Daten (siehe Bemerkung~\ref{rem:comp-protocol}).

\subsection{Benchmark: Nash-Frustration vs.\ funktionelle Reste}
\label{sec:frustration-benchmark}

Um explorativ zu beurteilen, ob Nash-Frustration Information über biologische Funktion tragen könnte, führten wir eine Benchmark-Analyse durch, die residuenspezifische Nash-Frustrationswerte mit experimentell annotierten funktionellen Resten für HP35 (1YRF) und Protein~G (1PGA) vergleicht. Für jedes Protein wurden Reste als \emph{funktionell} (katalytische Stellen, Bindungsschnittstellen, allosterische Regionen aus UniProt-Annotationen) oder \emph{passiv} (alle anderen) klassifiziert, und der Nash-Frustrationswert wurde als binärer Klassifikator mittels ROC- und Precision-Recall-Analyse bewertet. Tabelle~\ref{tab:frustration-benchmark} fasst die Ergebnisse zusammen.

\begin{table}[ht]
\centering
\caption{Explorativer Benchmark der Nash-Frustration als Klassifikator funktioneller Reste. AUC: Fläche unter der ROC-Kurve; AUPRC: Fläche unter der Precision-Recall-Kurve. Mittlere Frustrationswerte als Mittelwert $\pm$ Standardabweichung. Skript: \texttt{frustration\_benchmark.py}.}
\label{tab:frustration-benchmark}
\begin{tabularx}{\textwidth}{@{}L{0.18\textwidth} c c Y Y@{}}
\toprule
\textbf{Protein} & \textbf{AUC} & \textbf{AUPRC} & \textbf{Mittl.\ Frust.\ (funktionell)} & \textbf{Mittl.\ Frust.\ (passiv)} \\
\midrule
HP35 (1YRF) & 0{,}616 & 0{,}550 & $0{,}160 \pm 0{,}269$ & $0{,}031 \pm 0{,}082$ \\
Protein G (1PGA) & 0{,}688 & 0{,}618 & $0{,}229 \pm 0{,}247$ & $0{,}096 \pm 0{,}101$ \\
\bottomrule
\end{tabularx}
\end{table}

Die AUC-Werte von 0{,}62 (HP35) und 0{,}69 (Protein~G) zeigen eine schwache bis moderate explorative Trennung. Sie liegen in diesem kleinen Machbarkeitsnachweis über der Zufallsbasislinie, aber die aktuellen Benchmark-Dateien berichten noch keine Konfidenzintervalle und keinen formalen Signifikanztest für AUC-Werte über 0{,}5. Funktionelle Reste weisen in beiden Proteinen höhere mittlere Nash-Frustration auf als passive Reste ($0{,}160$ vs.\ $0{,}031$ für HP35; $0{,}229$ vs.\ $0{,}096$ für Protein~G), kompatibel mit, aber kein Beweis für die Interpretation, dass spieltheoretische Instabilität bevorzugt an biologisch aktiven Stellen lokalisiert ist. Die moderaten AUC-Werte sind angesichts des Charakters der Analyse als Machbarkeitsnachweis erwartbar: absolute Score-Schwellen bleiben unkalibriert, und das spieltheoretische Potential lässt Seitenketten-Wechselwirkungen, Lösungsmitteleffekte und Elektrostatik außer Acht. Diese Ergebnisse bleiben kompatibel mit der Interpretation, dass Nash-Frustration ein komplementäres Signal zur energetischen Frustration erfasst---eines, das Koordinationsversagen statt energetischer Suboptimalität reflektiert---aber weitere Kalibrierung erfordert, bevor es als eigenständiger Prädiktor dienen kann.

\begin{figure}[htbp]
\centering
\includegraphics[width=\textwidth]{benchmark/frustration_benchmark_roc.png}
\caption{Receiver-Operating-Characteristic (ROC)-Kurven für Nash-Frustration als binärer Klassifikator funktioneller versus passiver Reste. HP35 (1YRF, AUC\,=\,0{,}62) und Protein~G (1PGA, AUC\,=\,0{,}69) liegen beide über der Zufallsbasislinie (gestrichelte Diagonale), kompatibel mit der Hypothese einer bevorzugten Lokalisierung spieltheoretischer Instabilität an biologisch aktiven Stellen. Die moderaten AUC-Werte spiegeln den Charakter der Analyse als Machbarkeitsnachweis wider; Verbesserungen müssen durch reichere Potentiale und unabhängige Kontrollen gezeigt werden.}
\label{fig:frustration-benchmark-roc}
\end{figure}

\subsection{Genregulation als Signalspiel}

Genregulationsnetzwerke (GRNs) steuern Genexpressionsmuster und können als Signalspiele modelliert werden~\citep{GameBiochemistry}:

\begin{itemize}
    \item \textbf{Spieler:} Transkriptionsfaktoren und DNA-Bindungsstellen
    \item \textbf{Strategien:} Binden/Nicht-Binden (mit kontinuierlicher Affinität als Strategieintensität)
    \item \textbf{Payoff:} Organismische Fitness, vermittelt durch Genexpressionsniveaus
\end{itemize}

Die mathematische Verbindung zwischen Spieltheorie und Populationsdynamik wird durch die Replikatorgleichung hergestellt, die formal äquivalent zu Lotka--Volterra-Modellen ist~\citep{HofbauerSigmund1998}. Genauer gilt diese Äquivalenz für symmetrische Spiele mit linearen Fitnessfunktionen; bei frequenzabhängigen oder nichtlinearen Payoffs erfordert die Korrespondenz die Transformation $x_i = y_i / \sum_j y_j$ und beschränkt sich auf das Innere des Simplex (Hofbauer \& Sigmund 1998, Kapitel~7). Ein stabiler Fixpunkt in diesen Gleichungen entspricht einer ESS und somit einem Verhaltensattraktor, der die Langzeitstruktur des Systems bestimmt.

\paragraph{Illustratives Beispiel: das lac-Operon.}
Das Laktose-Verwertungssystem von \textit{E.\ coli} exemplifiziert die Signalspiel-Logik. Der lac-Repressor (LacI) und der CAP-Aktivator konkurrieren um die regulatorische Kontrolle des \textit{lac}-Operons. In Abwesenheit von Laktose stellt die LacI-Bindung ein Nash-Gleichgewicht dar: keine einseitige Änderung des Bindungszustands eines regulatorischen Faktors verbessert die organismische Fitness. Wenn Laktose vorhanden ist, inaktiviert Allolaktose LacI, verschiebt die Payoff-Matrix und etabliert ein neues Gleichgewicht mit Operon-Expression. Der bistabile Schalter zwischen diesen Zuständen ist präzise die Art von spieltheoretischem Attraktor, den der FST-III-Rahmen als ESS auf der Regulationsebene identifiziert.

% ===================================================================
\section{Das Free Energy Principle und Active Inference}
\label{sec:fep}
% ===================================================================

Das Free Energy Principle (FEP), primär entwickelt von Karl Friston, liefert eine vereinigende Theorie für biologische Selbstorganisation, die Informationstheorie und Thermodynamik verschmilzt~\citep{FristonFreeEnergy,FristonPDF}. \textbf{Hinweis:} Die Behauptung, dass FEP und MEPP kompatibel oder äquivalent sind, ist nicht unumstritten. FEP minimiert variationelle freie Energie (eine informationstheoretische Größe, die Surprise begrenzt), während MEPP die thermodynamische Entropieproduktionsrate maximiert. Die Verbindung zwischen diesen beiden Prinzipien ist ein aktives Diskussionsgebiet~\citep{FristonFreeEnergy}; das vorliegende Paper behandelt sie als komplementäre Perspektiven auf verschiedenen Analyseebenen, leitet aber nicht formal eines aus dem anderen ab. Für eine kritische Bewertung der philosophischen Ansprüche des FEP (z.\,B.\ Markovian Monism) gegenüber seiner wissenschaftlichen Reichweite als heuristisches Vereinheitlichungsprinzip siehe~\cite{SanchezCanizares2021}.

\subsection{Theoretischer Rahmen}

FEP postuliert, dass jedes biologische System, das sich erfolgreich gegen Entropie (Zerfall) verteidigt, seine ``variationelle freie Energie'' minimieren muss~\citep{FristonFreeEnergy}. Diese Größe repräsentiert eine mathematische obere Schranke der ``Surprise'' (Surprisal), die das System durch sensorische Inputs erlebt~\citep{FristonFreeEnergy}.

Ein biologischer Agent existiert innerhalb einer ``Markov-Decke''---einer statistischen Grenze, die interne Zustände von externen Zuständen trennt~\citep{FristonPDF}. Um zu überleben, muss der Agent sicherstellen, dass seine sensorischen Zustände innerhalb physiologischer Grenzen bleiben. Dies wird durch zwei Mechanismen erreicht:

\begin{enumerate}
    \item \textbf{Wahrnehmung:} Das System passt seine internen Modelle an, um die Ursachen sensorischer Daten besser vorherzusagen (Minimierung des Vorhersagefehlers)~\citep{FristonFreeEnergy}
    \item \textbf{Active Inference:} Das System handelt, um die Welt zu verändern oder sensorische Daten zu selektieren, die seinen internen Erwartungen entsprechen~\citep{FristonFreeEnergy}
\end{enumerate}

Durch Minimierung der freien Energie wird das System implizit zu einem Modell seiner Umgebung. Dieser Prozess ist eng mit MEPP verknüpft: während das System intern seine freie Energie minimiert (Ordnung schafft), maximiert es die Entropieproduktion in der Umgebung durch metabolische und agentive Aktivitäten~\citep{FristonPDF}.

In FST-III sind Verhaltensattraktoren diejenigen Zustände, in denen die freie Energie minimal ist und der Agent eine stabile, vorhersagbare Interaktion mit seiner Nische erreicht hat~\citep{FristonFreeEnergy}.

\subsection{Freie-Energie-Minimierung als Nash-Fixpunkt}

Die Verbindung zwischen FEP und Nash-Gleichgewicht kann unter bestimmten Bedingungen formalisiert werden. Sei $p(o,s_1,\dots,s_N)$ ein gemeinsames generatives Modell für $N$ interagierende Agenten, wobei jeder Agent $i$ eine faktorisierte approximative Posterior $q_i(s_i)$ über seine latenten Zustände kontrolliert. Die gemeinsame variationelle freie Energie ist:
\begin{equation}
F(q; o) = \mathbb{E}_q\!\left[\ln q(s) - \ln p(o, s)\right], \qquad q(s) = \prod_{i=1}^{N} q_i(s_i).
\end{equation}

Definiere ein Spiel, in dem jeder Agent $i$ die Kontrolle über $q_i$ hat und alle Agenten das Kostenfunktional $J_i(q_i, q_{-i}) := F(\prod_j q_j; o)$ teilen. Dann ist die beste Antwort von Agent $i$ gegen festes $q_{-i}$ präzise das koordinatenweise variationelle Update (Mean-Field-VI), und ein Nash-Gleichgewicht $q^* = (q_1^*, \dots, q_N^*)$ ist ein Fixpunkt verteilter Freie-Energie-Minimierung. Dies konstituiert ein \emph{Potentialspiel} mit Potential $F$ im Sinne von~\citet{MondererShapley1996}: Nash-Gleichgewichte fallen mit koordinatenweisen Minima der gemeinsamen freien Energie zusammen.

\paragraph{Gültigkeitsbereich der Brücke.} Diese formale Äquivalenz gilt für die Belief-Optimierungskomponente von FEP (variationelle Inferenz) und beruht auf der \emph{Potentialspiel}-Struktur: alle Agenten teilen dasselbe Kostenfunktional $F$. In der biologischen Realität haben interagierende Organismen typischerweise heterogene Zielfunktionen (Räuber vs.\ Beute, Wirt vs.\ Parasit), was die Potentialspiel-Annahme bricht. Die FEP-als-Nash-Brücke gilt daher strikt für intra-organismische Inferenz (wo alle Subsysteme das Ziel organismischen Überlebens teilen) und für symmetrische Mehr-Agenten-Szenarien; sie lässt sich nicht ohne Modifikation auf antagonistische ökologische Interaktionen übertragen. Für Active Inference, wo Agenten zusätzlich Policies $\pi_i$ selektieren, um die erwartete freie Energie $G(\pi_i, \pi_{-i})$ zu minimieren, wird die Nash-Bedingung: kein Agent kann einseitig seine Policy ändern, um seine erwartete freie Energie zu reduzieren---ein genuines Mehr-Agenten-Gleichgewicht, das Interaktion durch Zustandsdynamik involviert. Das vorliegende Paper entwickelt diese Erweiterung nicht formal; der Beitrag besteht darin, die präzise Demarkation zu liefern: FEP-als-Nash ist mehr als Metapher, wenn das Spiel explizit als variationelles/Potentialspiel über Verteilungen definiert ist, und bloß konzeptuell, wenn ``Nash'' als Synonym für ``Fixpunkt'' verwendet wird.

\subsection{Skalentrennung: Auflösung des FEP-MEPP-Paradoxons}
\label{sec:fep-mepp-paradox}

Eine persistente konzeptuelle Spannung in der theoretischen Biologie fragt, wie ein System seine interne freie Energie \emph{minimieren} (FEP) kann, während es gleichzeitig die Entropieproduktion in der Umgebung \emph{maximiert} (MEPP). Die folgende Konjektur adressiert diese Spannung durch Skalentrennung über die Markov-Decke.

\begin{conjecture}[FEP-MEPP-Kompatibilität via Markov-Decke]
\label{conj:fep-mepp}
Sei $\mathcal{B}$ ein selbstorganisierendes System mit Markov-Decke $\partial\mathcal{B}$, internen Zuständen $s$ und Umgebungszuständen $e$. Definiere:
\begin{itemize}
  \item $F_{\mathrm{int}}(q) :=
    \mathbb{E}_{q}\!\bigl[\ln q(s) - \ln p(o, s)\bigr]$,
    die variationelle freie Energie des internen Modells;
  \item $\dot{S}_{\mathrm{ext}} :=
    \frac{\dd S_{\mathrm{env}}}{\dd t} =
    \frac{1}{T}\,\frac{\dd Q}{\dd t}$,
    die Rate der Entropieproduktion in der Umgebung.
\end{itemize}
Wenn $\mathcal{B}$ $F_{\mathrm{int}}$ minimiert (FEP), indem es ein akkurates internes Modell aufrechterhält und Aktionen wählt, die Surprise reduzieren, dann:
\begin{enumerate}
  \item[\emph{(i)}] Die strukturelle Integrität von $\mathcal{B}$ wird aufrechterhalten: $\mathcal{B}$ widersteht der Auflösung und erhält die Markov-Decke (niedrige interne Entropie);
  \item[\emph{(ii)}] Die intakte, fernab-vom-Gleichgewicht-Struktur $\mathcal{B}$ fungiert als effiziente dissipative Maschine, die Energiegradienten aus der Umgebung durch metabolische und agentive Aktivitäten kanalisiert;
  \item[\emph{(iii)}] Unter allen Randbedingungen, die konsistent mit dem Überleben von $\mathcal{B}$ sind, tendiert die tatsächliche Dynamik zu Konfigurationen mit hoher (wenn auch nicht notwendigerweise maximaler) Rate der Entropieproduktion $\dot{S}_{\mathrm{ext}}$ in der Umgebung, konsistent mit der Stabilitätsinterpretation von MEPP (Hypothese~2 oben). Der Schlüsselmechanismus, der (ii) mit (iii) verbindet, ist \emph{kompetitive Selektion}: Unter mehreren intakten dissipativen Strukturen erwerben diejenigen, die Energiegradienten effizienter kanalisieren, mehr Ressourcen, replizieren schneller und verdrängen langsamere Dissipatoren. Dieser Selektionsdruck---abwesend bei passiven Strukturen wie Kristallen---treibt die Korrelation zwischen struktureller Integrität und hoher Entropieproduktion im biologischen Bereich.
\end{enumerate}
Kurz: \emph{interne Freie-Energie-Minimierung kann externe Entropieproduktions-Zunahme unterstützen}. Die beiden Prinzipien müssen nicht im direkten Widerspruch stehen, wenn sie komplementären Seiten der Markov-Decke zugeordnet werden.
\end{conjecture}

\noindent\textbf{Status:} Diese Konjektur ist gestützt auf qualitative Argumente. Ein formaler Beweis würde erfordern, dass FEP-minimierende Systeme notwendigerweise die Entropieproduktion in der Umgebung unter den genannten Randbedingungen maximieren -- ein Resultat, das in der Literatur nicht demonstriert wurde.

\begin{remark}[Spieltheoretische Interpretation]
In der Sprache des vorliegenden Papers entspricht die FEP-MEPP-Kompatibilität einem Zweistufen-Spiel:
\begin{itemize}
  \item \emph{Inneres Spiel} (innerhalb $\partial\mathcal{B}$): Agenten minimieren die variationelle freie Energie $F_{\mathrm{int}}$ über den Nash-Fixpunkt des in Abschnitt~\ref{sec:fep} beschriebenen Potentialspiels.
  \item \emph{Äußeres Spiel} (über $\partial\mathcal{B}$): Das Ensemble überlebender Organismen konkurriert um Energiegradienten. Organismen, die effizienter dissipieren (höheres $\dot{S}_{\mathrm{ext}}$), können langsamere Dissipatoren verdrängen, wodurch MEPP als makroskaliger Selektionskandidat unter den genannten Randbedingungen erscheint.
\end{itemize}
Dies formuliert das von Martyushev~\cite{Martyushev2010} und Sawada et al.~\citep{Toma2025} aufgeworfene MEPP-Paradoxon neu: MEPP gilt für das \emph{umgebungsgekoppelte} System, nicht für die internen Freiheitsgrade, wo minimale Entropieproduktion (Prigogine) lokal gelten kann. Die Nash-Frustration $\rho(J)$ aus Abschnitt~\ref{sec:molecular} wird hier als Modelldiagnostik für residuale Spannung zwischen innerer Optimalität und äußerem dissipativem Druck verwendet.
\end{remark}

\subsection{Nash-Gleichgewicht als Mean-Field-Game-Fixpunkt}\label{sec:mfg-fixpoint}

Die Potentialspiel-Charakterisierung der FEP-Nash-Brücke (Konjektur~\ref{conj:fep-mepp}) kann in den allgemeineren Rahmen der \emph{Mean Field Games} (MFG) eingebettet werden, wie sie von Lasry und Lions~\cite{LasryLions2007} eingeführt wurden. Man betrachte $2N$ Diederwinkel $\theta\in[-\pi,\pi]^{2N}$, die das Rückgrat eines faltenden Proteins parametrisieren. Jeder Rest wird als repräsentativer Agent behandelt, dessen Zustand durch seinen lokalen Torsionswinkel bestimmt wird. Im Kontinuumslimes (großes $N$) nimmt das gekoppelte Optimierungsproblem die Form eines MFG-Systems an, bestehend aus einer Hamilton--Jacobi-Bellman-(HJB-)Gleichung und einer Fokker--Planck-(FP-)Gleichung.

\textbf{Approximationsvorbehalt.} Die MFG-Formulierung nimmt an, dass jeder Agent mit der \emph{Populationsverteilung} der anderen Agenten interagiert, nicht mit einzelnen spezifischen Nachbarn. Für Proteine ist diese Mean-Field-Näherung am plausibelsten bei großen, dichten und annähernd homogenen Kontaktstrukturen. Reale Proteine besitzen jedoch spärliche, topologiespezifische Kontaktnetze; für kleine Proteine wie HP35 liefert die MFG-Einbettung daher eher ein formales Rahmenmodell als eine quantitativ genaue Beschreibung.

Die HJB-Gleichung beschreibt die optimale Steuerung jedes Agenten. Der zugehörige Hamiltonian lautet
\begin{equation}\label{eq:mfg-hamiltonian}
  H(\theta,p,m) \;=\; \tfrac{1}{2}\|p\|^{2}
    - \sum_{i<j}\int E_{ij}(\theta_i,\theta_j)\,m_j(\theta_j)\,\mathrm{d}\theta_j\,,
\end{equation}
wobei $p=\nabla V$ der konjugierte Impuls ist, $E_{ij}$ paarweise Rest-Rest-Wechselwirkungen kodiert und $m_j$ die Populationsdichte der Diederwinkel-Zustände an Position $j$ ist. Die Fokker--Planck-Gleichung propagiert die Agentendichte unter dem optimalen Drift und thermischem Rauschen:
\begin{equation}\label{eq:mfg-fp}
  \partial_t m \;=\; \nabla\!\cdot\!\bigl(m\,\nabla V\bigr) + \sigma\,\Delta m\,,
  \qquad \sigma = k_B T / \gamma\,.
\end{equation}
Ein MFG-Nash-Gleichgewicht ist ein Paar $(V^{*},m^{*})$, das beide Gleichungen simultan erfüllt. Im thermischen Gleichgewicht ist die stationäre Dichte
\begin{equation}\label{eq:mfg-boltzmann}
  m^{*}(\theta) \;\propto\; \exp\!\bigl(-V(\theta)/k_B T\bigr)\,,
\end{equation}
was identisch zur Boltzmann-Verteilung aus der Langevin-Fokker--Planck-Beschreibung der Proteinfaltung ist. Das bestehende Potentialspiel-Argument (Abschnitt~\ref{sec:fep}) wird als Spezialfall wiedererhalten: Im \emph{statischen Limes} $\partial_t=0$ reduziert sich die HJB-Gleichung auf $\nabla V=0$ (Agenten sitzen an Optima), und die Fokker--Planck-Gleichung wird zu $\nabla\!\cdot(m\nabla V)=0$, was automatisch erfüllt ist. Das resultierende System ist ein klassisches Potentialspiel mit gemeinsamen Kostenfunktional $V$.

Die entscheidende Existenzbedingung für MFG-Nash-Gleichgewichte ist die \textbf{Lasry--Lions-Monotonie}~\cite{LasryLions2007}: Die Interaktionskosten müssen monoton im Populationsmaß sein, d.\,h.\ die Abbildung $m\mapsto -\sum_{i<j}\int E_{ij}\,m_j\,\mathrm{d}\theta_j$ muss monoton sein. In idealisierten Fällen mit dominierenden \emph{repulsiven} Wechselwirkungen (Excluded-Volume-Effekte, sterische Konflikte) ist diese Bedingung plausibel, da eine Erhöhung der Dichte bei einer gegebenen Konfiguration die Kosten für alle Agenten erhöht; in dieser Darstellung würde dies Eindeutigkeit und Konvergenz zu einem Nash-Gleichgewicht stützen. Für \emph{stark attraktive} Wechselwirkungen hingegen ist die Monotonie verletzt: Kooperative Bindung senkt die Kosten, wenn mehr Agenten aggregieren, was mehrere selbstverstärkende Fixpunkte erzeugt und Konvergenzgarantien zerstört.

Die physikalische Interpretation ist nützlich, aber modellabhängig: Lasry--Lions-Monotonie entspricht in dieser Darstellung der Abwesenheit stark kooperativer Fehlfaltung. Wenn Monotonie gilt, besitzt das idealisierte MFG-System Konvergenzgarantien zu einem stabilen Nash-Gleichgewichtskandidaten für die native Faltung. Wenn sie verletzt ist, kann das System mehrere konkurrierende Fixpunkte zulassen, die fehlgefalteten Aggregaten entsprechen.

\begin{remark}[Versagensmodi]\label{rem:mfg-failure}
Die idealisierte MFG-Darstellung würde Nicht-Konvergenz in Systemen mit stark kooperativer Fehlfaltung erwarten lassen (Prionen, Amyloid-Fibrillen), in denen die Lasry--Lions-Monotonie verletzt ist. Dies ist eine nützliche Geltungsbereichsbedingung statt eine Schwäche des Rahmens: solche Systeme liegen außerhalb der einfachsten Stable-Nash-Domäne. Prionenpropagation und Amyloid-Nukleation sind plausible Beispiele für Regime, in denen die Potentialspiel-Struktur zusammenbricht und das Multi-Agenten-System in kinetisch stabilisierten, thermodynamisch metastabilen Zuständen gefangen wird. Die MFG-Einbettung schärft somit die Gültigkeitsaussage der FEP-Nash-Brücke: Sie gilt für Faltungssysteme, die Monotonie erfüllen, während ihre Versagensmodi mit bekannten Pathologien der Proteinaggregation verglichen werden können.
\end{remark}

% ===================================================================
\section{Boolesche Netzwerke und der Rand des Chaos}
\label{sec:boolean}
% ===================================================================

Auf der Ebene der Genregulation manifestiert sich Stabilität in der Dynamik Boolescher Netzwerke, wie sie von Stuart Kauffman eingeführt wurden~\citep{BooleanPost,KauffmanOrigins}. Boolesche Netzwerkmodelle sind ein aktives Forschungsfeld; für umfassende Behandlungen von Netzwerkmotiven und Designprinzipien siehe~\citet{AlonSystems2007}, und für Robustheit und Evolvabilität in genetischen Netzwerken siehe~\citet{WagnerRobustness2005}.

\subsection{Netzwerkdynamik}

Diese Netzwerke modellieren Gene als binäre Schalter, deren Zustände durch regulatorische Logik verknüpft sind~\citep{BooleanPost}. Die Langzeitdynamik solcher Netzwerke fließt in Attraktoren, die biologisch als spezifische Zelltypen oder funktionelle Zustände interpretiert werden~\citep{BooleanEvolution}.

Kauffman argumentierte, dass Netzwerkstabilität von der Konnektivität und dem Typ der Booleschen Funktionen abhängt~\citep{BooleanPost}. Netzwerke am ``Rand des Chaos''---einer kritischen Übergangsregion zwischen rigider Ordnung und unvorhersagbarem Chaos---weisen hypothetisch hohe Evolvabilität und Informationsverarbeitungskapazität auf, obwohl diese Behauptung in der Literatur umstritten bleibt~\citep{KauffmanOrigins}:

\begin{table}[htbp]
\centering
\caption{Regime Boolescher Netzwerke und häufig vorgeschlagene biologische Interpretationen. Für kritische Netzwerke wird eine Balance von Robustheit und Anpassungsfähigkeit hypothetisiert; weder maximale Evolvabilität noch ein universelles biologisches Optimum werden hier vorausgesetzt.}
\label{tab:network-regimes}
\begin{tabularx}{\textwidth}{@{}L{0.22\textwidth} L{0.28\textwidth} Y@{}}
\toprule
\textbf{Netzwerkregime} & \textbf{Dynamisches Verhalten} & \textbf{Biologische Funktion} \\
\midrule
Geordnet & Eingefrorene Zustände, geringe Variabilität & Strukturelle Integrität, Basalmetabolismus \\
Kritisch (Rand des Chaos) & Komplexe Muster, hoher Informationsfluss & Adaptives Verhalten, Differenzierung, Lernen \\
Chaotisch & Hohe Sensitivität, Instabilität & Pathologische Zustände, Malignität \\
\bottomrule
\end{tabularx}
\end{table}

Diese Attraktoren in Genregulationsnetzwerken bilden die materielle Basis für die spieltheoretischen Strategien auf höheren Ebenen. Eine ESS auf der Makroebene erfordert stabile genetische Attraktoren auf der Mikroebene, die zuverlässig die notwendigen phänotypischen Merkmale produzieren~\citep{BooleanMicroarray,KauffmanOrigins}.

\paragraph{Boolesche Attraktoren als Nash-Gleichgewichte.}
Die Verbindung zwischen Booleschen Netzwerk-Attraktoren und Nash-Gleichgewichten kann explizit gemacht werden. Jedes Gen $g_i$ in einem Booleschen Netzwerk von $n$ Genen ist ein ``Spieler'', dessen Strategie sein Expressionszustand $\sigma_i \in \{0, 1\}$ ist. Die Boolesche Update-Funktion $f_i(\sigma_{-i})$ definiert die beste Antwort von Gen $i$ gegeben die Zustände seiner Regulatoren. Ein Fixpunkt-Attraktor $\sigma^* = (f_1(\sigma^*_{-1}), \dots, f_n(\sigma^*_{-n}))$ ist dann präzise ein Nash-Gleichgewicht in reinen Strategien dieses $n$-Spieler-Spiels: kein Gen kann einseitig seinen Zustand ändern, um seinen ``Payoff'' (definiert als Übereinstimmung mit der regulatorischen Logik) zu verbessern. \textbf{Hinweis:} Diese Nash-Charakterisierung ist formal statt substantiell---sie reformuliert die Fixpunktbedingung in spieltheoretischer Sprache. Der nicht-triviale Gehalt entsteht erst, wenn der Payoff an ein externes Fitnessmaß gekoppelt wird (z.\,B.\ organismische Lebensfähigkeit), sodass nicht alle Booleschen Fixpunkte gleichermaßen ``Nash-stabil'' unter Selektion sind. Grenzzyklus-Attraktoren entsprechen periodischen Orbits statt reinen Nash-Gleichgewichten, können aber als korrelierte Gleichgewichte im wiederholten Spiel interpretiert werden. Diese spieltheoretische Perspektive auf Boolesche Netzwerke liefert die formale Brücke zwischen Kauffmans Attraktor-Dynamik und der hierarchischen Nash-Architektur von FST-III.

% ===================================================================
\section{Zelluläre Ebene: Multizellularität als Koalitionsspiel}
\label{sec:multicellularity}
% ===================================================================

\subsection{Das Kooperationsproblem}

In einem multizellulären Organismus ``opfern'' einzelne Zellen ihr Reproduktionspotential zugunsten des Ganzen. Somatische Zellen reproduzieren sich nicht direkt; nur Keimbahnzellen übertragen Gene.

Dies repräsentiert extreme Kooperation: Zellen adoptieren eine Strategie (Differenzierung, Nicht-Reproduktion), die dem Kollektiv auf individuelle Kosten nützt.

\subsection{Koalitionsspiel-Formulierung}

\begin{itemize}
    \item \textbf{Spieler:} Einzelne Zellen innerhalb des Organismus
    \item \textbf{Strategien:} Kooperieren (Entwicklungsprogramm folgen) oder Defektieren (autonom proliferieren)
    \item \textbf{Payoff:} Genetische Fitness (über Keimbahn übertragen)
\end{itemize}

In erster Näherung sind alle Zellen in einem Organismus genetisch nahezu identisch (klonal) und teilen die überwältigende Mehrheit ihres Genoms. Daher entspricht der Payoff für jede Zelle näherungsweise dem Payoff der Keimbahn: wenn sich der Organismus reproduziert, werden die gemeinsamen Gene übertragen. (In der Praxis führen somatische Mutationen, V(D)J-Rekombination in Immunzellen und mitochondriale Heteroplasie zu begrenzter genetischer Divergenz, was präzise das Krebs-als-Defektion-Szenario ermöglicht, das unten diskutiert wird.)

\begin{proposition}[Multizellularität als Nash-Gleichgewicht]
\label{prop:multicellularity}
In einem klonalen Organismus, in dem alle Zellen den Verwandtschaftsgrad $r \approx 1$ teilen, ist Kooperation (Befolgen des Entwicklungsprogramms) ein Nash-Gleichgewicht, sofern die individuellen Kosten der Kooperation geringer sind als der inklusive Fitnessnutzen: $c < r \cdot b$ (Hamiltons Regel~\citep{Hamilton1964}). Im Grenzfall $r = 1$ genügt jeder positive Nettonutzen $b > c$. Der Beitrag hier besteht in der Einbettung dieses klassischen Resultats in den hierarchischen FST-Rahmen; das Nash-Gleichgewicht wird aufrechterhalten, solange die klonale Identität erhalten bleibt.
\end{proposition}

\subsection{Krebs als Koalitionszusammenbruch}

Krebs als Zusammenbruch multizellulärer Kooperation ist ein gut etablierter Rahmen~\citep{AktipisCancer}. Die spieltheoretische Rahmung hier folgt Aktipis et al.\ (2015), die evolutionäre Spieltheorie anwenden, um somatische Evolution und Krebsprogression zu erklären. Der FST-III-Beitrag besteht darin, diesen Rahmen in das hierarchische Nash-Gleichgewichts-Bild einzubetten: Krebs repräsentiert nicht bloß einen selektiven Vorteil für einzelne Zellen, sondern einen Phasenübergang auf der Koalitionsebene, analog zum Parasitenproblem in der präbiotischen Chemie, das in FST-II diskutiert wird. Formal repräsentiert Krebs einen Koalitionszusammenbruch:
\begin{enumerate}
    \item Somatische Mutationen erzeugen genetische Divergenz
    \item Mutierte Zellen teilen nicht mehr alle Gene mit anderen Zellen
    \item Defektion (Proliferation) nützt nun der mutierten Linie
    \item Das Nash-Gleichgewicht bricht zusammen; Krebszellen ``defektieren''
\end{enumerate}

Dies verbindet sich mit dem Parasitenproblem in Hyperzyklen~\citep{FSTII}: sowohl Krebs als auch chemische Parasiten nutzen kooperative Systeme aus, indem sie Bedingungen brechen, die Kooperation stabilisieren. Die Parallele zwischen Krebs als Koalitionszusammenbruch und dem Hyperzyklus-Parasitenproblem wird als \emph{strukturelle Analogie} präsentiert, nicht als abgeleitetes Resultat. Eine formale Verbindung würde erfordern zu zeigen, dass onkogene Mutationen auf parasitäre Replikatoren im Sinne von Eigen \& Schuster abbilden; dies bleibt ein offenes Problem.

% ===================================================================
\section{Skalierung und die Renormierungsgruppe}
\label{sec:renormalization}
% ===================================================================

Ein zentrales Element von FST-III ist die Erkenntnis, dass biologische Prinzipien skaleninvariant über viele Ebenen operieren~\citep{ScaleRGFrontiers,ScaleRGResearch}.

\subsection{Der RG-Rahmen}

\textbf{Terminologischer Vorbehalt:} Der Begriff ``Renormierungsgruppe'' wird in diesem Abschnitt in einem erweiterten, metaphorischen Sinne verwendet. In der Physik der kondensierten Materie bezieht sich die Renormierungsgruppe (RG) auf einen präzisen mathematischen Rahmen mit Flussgleichungen, Fixpunkten und Skalenexponenten~\citep{ScaleRGResearch}. Die vorliegende Verwendung---die beschreibt, wie effektive spieltheoretische Parameter sich ändern, wenn man auf sukzessive größeren biologischen Skalen betrachtet---ruft die \emph{konzeptuelle Logik} der Grobkörnung auf, liefert aber keine RG-Flussgleichungen, Fixpunktanalyse oder dimensionelle Skalengesetze für biologische Systeme. Ob eine rigorose RG-Behandlung evolutionärer Übergänge erreichbar ist, bleibt eine offene Frage~\citep{ScaleRGFrontiers}.

Mit diesem Vorbehalt im Hinterkopf: Die Grobkörnungsperspektive postuliert, dass biologische Systeme oft hierarchisch organisiert sind---von fraktaler Verzweigung in Blutgefäßen bis zu hierarchischen neuronalen Architekturen~\citep{SelfSimilarNetworks}.

Renormalizing Generative Models (RGM) zeigen, dass Prinzipien der freien Energie und spieltheoretischen Optimierung rekursiv auf jeder Ebene angewandt werden~\citep{ScaleRGFrontiers}. Ein Organismus kann als Grobkörnung über Milliarden von Zellen betrachtet werden, wobei Zellen als ``Atome'' in einem Spiel höherer Ordnung fungieren, dessen Regeln durch die evolutionäre Geschichte der Spezies festgelegt sind~\citep{ReplicatorOptimization}.

\subsection{Das skalenrelative Kernmodell}

Diese Instanziierung wird durch skalenrelative Fitnessmodelle der Mehrebenenselektion beschrieben~\citep{Okasha2006,Michod1999,ReplicatorOptimization}:

\begin{table}[htbp]
\centering
\caption{Skalenrelatives Kernmodell: Fitnessdefinitionen und Strategietypen auf jeder biologischen Organisationsebene.}
\label{tab:scale-kernel}
\begin{tabular}{ll}
\toprule
\textbf{Ebene} & \textbf{Fitness/Strategie} \\
\midrule
Zellulär & Replikationsrate als Fitness, Mutationen als Strategiewechsel \\
Organismisch & Darwinsche Fitness, Verhaltensstrategien \\
Agentiv & Intentionalität, Lernen durch Imitation \\
Institutionell & Legitimität und soziale Kooperation als ESS \\
\bottomrule
\end{tabular}
\end{table}

Diese hierarchische Struktur stellt sicher, dass Konflikte zwischen Skalen (z.\,B.\ Krebs als Defektion einzelner Zellen gegen den Organismus) durch übergeordnete Selektionsdrücke minimiert werden~\citep{ReplicatorOptimization}.

\subsection{Hauptübergänge als RG-Schritte}

Die Hauptübergänge der Evolution, identifiziert von Szathm\'ary und Maynard Smith~\citep{SzathmaryMaynardSmith1995}, bilden auf diesen Rahmen ab:

\begin{itemize}
    \item Gene $\to$ Chromosomen (Block 1)
    \item Prokaryoten $\to$ Eukaryoten (Block 2)
    \item Einzeller $\to$ Vielzeller (Block 3)
    \item Individuen $\to$ Gesellschaften (Block 4)
\end{itemize}

Jeder Übergang schafft eine neue Ebene von ``Spielern'' aus Koalitionen niedrigerer Spieler und erhält spieltheoretische Struktur durch Grobkörnung.

\paragraph{Ein Toy-Renormierungsschritt für räumliche evolutionäre Spiele.}
Betrachte ein 2D-Gitter-Spiel mit Strategien $\sigma_x\in\{C,D\}$ und Payoff-Matrix
$A=\begin{pmatrix}R&S\\ T&P\end{pmatrix}$.
Für einen $2\times 2$-Block $B$ definiere eine Blockstrategie $\Sigma_B=\mathcal B(\sigma|_B)$ durch Mehrheitsregel.
Für benachbarte Blöcke $B,B'$ definiere den effektiven Block-Payoff durch Randmittelung

\[
A'_{s,t}
:=\frac{1}{|\partial(B,B')|}\,
\mathbb E\!\left[\sum_{\langle x,y\rangle\in\partial(B,B')} A_{\sigma_x,\sigma_y}\ \Big|\ \Sigma_B=s,\ \Sigma_{B'}=t\right],
\qquad s,t\in\{C,D\}.
\]

In einer Frozen-Interior-Approximation, sei
$q_s=\mathbb P(\sigma_x=C\mid \Sigma_B=s,\ x\in\partial B)$.
Dann erfüllt die RG-Abbildung $R_2:(A,\beta)\mapsto(A',\beta')$ die explizite Mischungsformel

\[
A'_{s,t}
= q_s q_t\,R + q_s(1-q_t)\,S + (1-q_s)q_t\,T + (1-q_s)(1-q_t)\,P,
\]

mit $q_C\approx 1-\varepsilon$ und $q_D\approx \varepsilon$, kontrolliert durch die Within-Block-Dynamik (Selektionsstärke $\beta$ und interne Durchsetzungsmechanismen).
Fixpunkte umfassen die geordneten Phasen (All-$C$, All-$D$), während Änderungen in Within-Block-Einschränkungen den RG-Fluss durch Modifikation von $q_C-q_D$ verschieben.
Dies liefert einen konkreten Sinn, in dem evolutionäre Hauptübergänge Grobkörnungsschritten zusammen mit der Aktivierung neuer relevanter Operatoren (z.\,B.\ Policing, Bottlenecks) entsprechen, die kooperative Makrovariablen stabilisieren. Biologisch entspricht die Blockstrategie $\Sigma_B$ einem höherstufigen Phänotyp (z.\,B.\ multizelluläre Kooperation), während die Randparameter $q_C, q_D$ Durchsetzungsmechanismen kodieren (z.\,B.\ Immunüberwachung, reproduktive Bottlenecks), die Innergruppen-Defektion unterdrücken.

% ===================================================================
\section{Ökosystem-Ebene: Das ökologische Spiel}
\label{sec:ecosystem}
% ===================================================================

\subsection{Arten als Spieler}

Im ökologischen Spiel:
\begin{itemize}
    \item \textbf{Spieler:} Artenpopulationen
    \item \textbf{Strategien:} Ökologische Rollen (Räuber, Beute, Konkurrent, Mutualist)
    \item \textbf{Payoff:} Populationswachstumsrate
\end{itemize}

\subsection{Ökologische stationäre Zustände als mögliche spieltheoretische Analogien}

Eine ökologische Nische ist nicht selbst ein Nash-Gleichgewicht; sie beschreibt Umweltanforderungen und Interaktionen einer Art. Eine Nash- oder ESS-Deutung ist erst nach expliziter Festlegung vererbbarer Strategien, frequenzabhängiger Payoffs und zulässiger einseitiger Abweichungen möglich. Populationsdynamische Fixpunkte können dann unter einer expliziten Replikator--Lotka--Volterra-Abbildung mit Spielgleichgewichten verglichen werden; ihre dynamische Stabilität ist jedoch separat nachzuweisen.

\paragraph{Beispiel: Räuber-Beute-Koexistenz.}
In einem klassischen Lotka--Volterra-Räuber-Beute-System erfüllt der innere Fixpunkt---an dem Räuber- und Beutepopulationen bei konstanten Dichten koexistieren---die Nash-Gleichgewichtsbedingung des ökologischen Spiels: weder die Beuteart (durch Änderung ihrer Reproduktionsstrategie) noch der Räuber (durch Änderung seiner Jagdstrategie) kann an diesem Punkt einseitig seine Populationswachstumsrate verbessern. \textbf{Vorbehalt:} Im einfachen Lotka--Volterra-Modell ist dieser Fixpunkt ein Zentrum (neutral stabil mit geschlossenen Orbits), kein asymptotisch stabiler Attraktor; ESS-artige Stabilität erfordert zusätzliche Mechanismen wie dichteabhängige Mortalität oder räumliche Struktur. Die Äquivalenz zwischen Lotka--Volterra-Fixpunkten und Replikatordynamik-Gleichgewichten~\citep{HofbauerSigmund1998} bildet das ökologische Spiel auf die evolutionäre Dynamik ab, aber der Stabilitätstyp des resultierenden Gleichgewichts hängt von der spezifischen Payoff-Struktur ab.

\subsection{Massenaussterben als Phasenübergänge}

Massenaussterbe-Ereignisse repräsentieren Phasenübergänge im ökologischen Spiel:
\begin{enumerate}
    \item Umweltstörung verändert die Payoff-Matrix
    \item Vorherige Nash-Gleichgewichte werden instabil
    \item System geht in neues Gleichgewicht mit anderer Artenzusammensetzung über
\end{enumerate}

Auf struktureller Ebene gleicht diese Logik Phasenübergängen in physikalischen Systemen allgemein: veränderte Randbedingungen destabilisieren bestehende Gleichgewichte und treiben Übergänge zu neuen Konfigurationen. Es sei angemerkt, dass diese spieltheoretische Rahmung ein \emph{post hoc} interpretatives Modell ist: Massenaussterben sind komplexe Ereignisse, angetrieben durch Asteroideneinschläge, Vulkanismus oder Ozean-Anoxie, und inwieweit sie produktiv als Payoff-Matrix-Perturbationen modelliert werden können---statt als katastrophale Regimewechsel außerhalb des Gleichgewichtsrahmens---bleibt eine empirische Frage~\citep{GouldEldredge1977}.

% ===================================================================
\section{Synthese: Drei Phasen biologischer Organisation}
\label{sec:synthesis}
% ===================================================================

FST-III liefert ein integriertes Bild biologischer Organisation, das sich in drei Phasen entfaltet. Tabelle~\ref{tab:biological-game-hierarchy} fasst die hierarchische Spiel- und Stabilitätsarchitektur über alle Organisationsebenen zusammen.
\begin{table}[htbp]
\centering
\caption{Synthesematrix: Hierarchische Spiel- und Stabilitätsarchitektur über biologische Organisationsebenen. Jede Ebene instanziiert spieltheoretische Gleichgewichte unter thermodynamischen und dissipativen Nebenbedingungen.}
\label{tab:biological-game-hierarchy}
\small
\begin{tabularx}{\textwidth}{@{}L{0.16\textwidth} L{0.18\textwidth} L{0.20\textwidth} L{0.22\textwidth} Y@{}}
\toprule
\textbf{Ebene} & \textbf{Spieler} & \textbf{Strategieraum} & \textbf{Auszahlung / Kosten} & \textbf{Gleichgewicht \& Stabilität} \\
\midrule
Molekular & Aminosäurereste & Diederwinkel $(\phi, \psi)$ & $-\Delta G$ (Freie Energie) & Nash-Konformation / Resolventen-Stabilität \\
Genregulatorisch & Transkriptionsfaktoren / DNA & Bindungszustände $\sigma_i \in \{0,1\}$ & Fitness / Expressionslogik & Boolesche Fixpunkte / Signal-Attraktoren \\
Zellulär & Somatische / Keimbahn-Zellen & Kooperation vs.\ Defektion & Inklusive Fitness ($c < r \cdot b$) & Klonaler Nash-Zustand (Krebs = Kollaps) \\
Agentisch & Organismus / Subsysteme & Active-Inference-Policies $\pi$ & Variationelle Freie Energie $F$ & Dynamischer Bayesscher Nash-Attraktor \\
Ökologisch & Artenpopulationen & Nischenrollen / Strategien & Populationswachstumsrate & Lotka--Volterra ESS / Resolventen-Resilienz \\
\bottomrule
\end{tabularx}
\end{table}


\subsection{Phase 1: Thermodynamische Dissipation}

Der Energiefluss des jungen Universums zwingt einfache Materie in dissipative Strukturen. Hier wirkt MEPP als Stabilitätsfilter (Abschnitt~\ref{sec:thermodynamics}): unter den vielen möglichen Konfigurationen persistieren diejenigen mit resolvent-stabilen Dissipationsmustern. Es ist die Ära schneller Entropieproduktion und der Entstehung von Englands dissipativen Adaptoren~\citep{EcosystemMEPP}.

\subsection{Phase 2: Strategische Konsolidierung}

Sobald dissipative Strukturen die Schwelle zur Selbstreplikation überschreiten, führt Wettbewerb um Ressourcen zu spieltheoretischer Dynamik. Die Entstehung von Evolutionarily Stable Strategies (ESS) stabilisiert molekulare Populationen gegen parasitäre Invasion (Abschnitt~\ref{sec:molecular}). Die ``Spieler'' sind noch rein chemisch, aber ihre Interaktionen können bereits eine Nash-Gleichgewichts-Struktur aufweisen, die in späteren Ebenen wiederkehrt~\citep{MaynardSmithPrice1973}. Die evolutionären Hauptübergänge (Abschnitt~\ref{sec:renormalization}) körnen niedrigere Spielerebenen iterativ zu höheren Koalitionen, wobei die spieltheoretische Logik auf neuen Skalen teilweise erhalten bleiben kann.

\subsection{Phase 3: Kognitive Inferenz}

Organismen entwickeln interne Modelle ihrer Umgebung, um ihre freie Energie zu minimieren. Die Entstehung der Markov-Decke (Abschnitt~\ref{sec:fep}) erlaubt die Trennung von Selbst und Welt und motiviert die FEP-MEPP-Kompatibilität der Konjektur~\ref{conj:fep-mepp}: interner Modellbau (Freie-Energie-Minimierung) könnte äußere dissipative Effizienz (hohe Entropieproduktion) unterstützen.

Diese dritte Phase markiert einen qualitativen Wandel in der spieltheoretischen Hierarchie. In Phase~2 sind die ``Spieler'' molekulare oder zelluläre Agenten, deren Strategien durch physikalische Chemie bestimmt werden. In Phase~3 erwerben Agenten \emph{generative Modelle} ihrer Umgebung: die Strategien sind nicht mehr fixierte phänotypische Zustände, sondern Policies---Sequenzen von Aktionen, bedingt auf inferierte Umgebungszustände. Active Inference~\citep{FristonFreeEnergy} liefert den Mechanismus: Agenten wählen Aktionen, die die erwartete freie Energie minimieren (d.\,h.\ sie handeln, um ihre Vorhersagen zu bestätigen oder Unsicherheit zu reduzieren). Dies konvertiert die statischen Nash-Gleichgewichte von Phase~2 in dynamische Bayessche Nash-Gleichgewichte, bei denen die Strategie jedes Agenten von seinen Überzeugungen über andere abhängt. Die verschachtelte Hierarchie der Inferenz---von interozeptiver Homöostase über Wahrnehmung bis zur Planung---spiegelt die Grobkörnungslogik von Abschnitt~\ref{sec:renormalization}: jede Inferenzebene integriert über die Freiheitsgrade darunter und produziert ein zunehmend abstraktes Spiel auf höheren zeitlichen und räumlichen Skalen. Der Endpunkt dieser Hierarchie---Bewusstsein, Sprache und Kultur---kann als hochkomprimierte Form der spieltheoretischen Stabilitätslogik \emph{interpretiert} werden, wobei die ``Strategien'' symbolische Repräsentationen und der ``Payoff'' prädiktive Genauigkeit über beliebig lange Zeithorizonte ist~\citep{FristonPDF}. Diese Extrapolation ist spekulativ und geht über das hinaus, was der vorliegende Rahmen formal ableiten kann.

% ===================================================================
\section{Testbare Vorhersagen}
\label{sec:predictions}
% ===================================================================

\noindent\textit{Hinweis:} Die folgenden Vorhersagen sind mit P1--P5 gekennzeichnet, um sie von den mathematischen Propositionen (z.\,B.\ Proposition~\ref{prop:multicellularity}) in den vorangehenden Abschnitten zu unterscheiden. Propositionen sind formale Aussagen innerhalb des theoretischen Rahmens; Vorhersagen (Predictions) sind empirisch testbare Hypothesen, die daraus abgeleitet werden.

\subsection{P1: Proteinfaltungslandschaften}

\begin{hypothesis}[Nash-Frustration sagt faltungsrelevante Reste vorher]
Die Energielandschaft der Proteinfaltung kann durch Nash-Gleichgewichts-Stabilitätsanalyse charakterisiert werden. Reste mit hoher Nash-Frustration entsprechen experimentell identifizierten Faltungskeimen und allosterischen Stellen.
\end{hypothesis}

\textbf{Test:} Berechne Nash-Frustrationswerte für Proteine mit experimentell charakterisierten Faltungsintermediaten (z.\,B.\ via Wasserstoff-Deuterium-Austausch oder $\phi$-Wert-Analyse). Ein Korrelationskoeffizient $r > 0{,}3$ zwischen Nash-Frustrations-Rang und experimentellen $\phi$-Werten würde positive Evidenz darstellen.

\subsection{P2: Genregulationsnetzwerk-Attraktoren}

\begin{hypothesis}[GRN-Attraktoren als Nash-Gleichgewichte]
Zelltyp-Attraktoren in Genregulationsnetzwerken können durch Nash-Analyse des TF-DNA-Bindungsspiels modelliert und getestet werden.
\end{hypothesis}

\textbf{Test:} Berechne für gut charakterisierte GRNs Nash-Gleichgewichte des TF-DNA-Bindungsspiels, etwa für das \textit{Drosophila}-Segmentpolaritätsnetzwerk oder für ein hämatopoietisches Differenzierungsnetzwerk von Säugetieren. Vergleiche Anzahl und Identität der vorhergesagten Attraktorzustände mit experimentell beobachteten Zelltypen. Erfolgskriterium: Die vorhergesagte Attraktorzahl sollte mit der bekannten Zelltypzahl innerhalb eines Faktors zwei übereinstimmen.

\subsection{P3: Krebs und Koalitionsstabilität}

\begin{hypothesis}[Krebs als Koalitionszusammenbruch]
Die Krebsinzidenz korreliert invers mit der multizellulären Koalitionsstabilität, gemessen an der spieltheoretischen Defektionsschwelle.
\end{hypothesis}

\textbf{Test:} Über Gewebetypen hinweg, berechne die Defektionsschwelle $\Delta r$, bei der Hamiltons Bedingung $c < r \cdot b$ versagt (gegeben gewebespezifische Mutationsraten und Kooperationskosten). Sage vorher, dass Gewebe mit niedrigeren $\Delta r$-Schwellen höhere Krebsinzidenz aufweisen. Vergleiche mit epidemiologischen Daten aus der SEER-Datenbank. Ein besonders stringenter Test ergibt sich aus Petos Paradoxon: große, langlebige Arten zeigen keine proportional höheren Krebsraten trotz mehr Zellen und Zellteilungen. Das Koalitionsmodell erwartet, dass diese Arten stärkere Durchsetzungsmechanismen aufweisen könnten (niedrigeres effektives $\Delta r$), eine Vorhersage, die durch Vergleich der Immunüberwachungsintensität über Arten verschiedener Körpergrößen testbar ist.

\subsection{P4: Metabolische Rate und Fortpflanzungserfolg}

\begin{hypothesis}[MEPP-Fitness-Korrelation]
Über Arten hinweg korreliert die massenspezifische Stoffwechselrate mit der Fortpflanzungsrate nach Kontrolle für Körpergröße.
\end{hypothesis}

\textbf{Test:} Stelle Stoffwechselraten und Fortpflanzungsleistungen über Taxa zusammen. Eine positive Partialkorrelation (kontrolliert für Körpermasse) unterstützt die MEPP-Fitness-Verbindung. Kleibers Gesetz liefert die Körpergrößen-Skalierung. \textbf{Kontext:} Die Metabolic Theory of Ecology~\citep{BrownMTE2004} etabliert bereits quantitative Beziehungen zwischen Stoffwechselrate, Körpergröße und Life-History-Merkmalen. Die FST-III-Vorhersage geht über MTE hinaus, indem sie die Korrelation nicht bloß auf allometrische Skalierung, sondern auf den thermodynamischen Stabilitätsfilter (MEPP) zurückführt: Arten, die pro Masseeinheit effizienter dissipieren, sind nicht nur schnellere Reproduzierer, sondern sollten auch resilienter gegenüber Störungen sein.

\subsection{P5: Ökosystem-Resilienz}

\begin{hypothesis}[Ökosystem-Resilienz aus Nash-Stabilität]
Ökosystem-Resilienz wird als korreliert mit der spektralen Lücke der Nash-Stabilitätsmatrix des ökologischen Gleichgewichts hypothesiert.
\end{hypothesis}

\textbf{Test:} Berechne für gut charakterisierte Ökosysteme die Jacobi-Eigenwert-Lücke am Nash-Gleichgewicht, etwa für das Yellowstone-Nahrungsnetz oder für Korallenriff-Gemeinschaften. Sage vorher, dass Ökosysteme mit größerer spektraler Lücke sich schneller von Störungen erholen. Vergleiche dies mit empirischen Resilienzmaßen wie der Erholungszeit nach Störungsereignissen.

% ===================================================================
\section{Diskussion}
\label{sec:discussion}
% ===================================================================

\subsection{Was ist wirklich neu?}

Die klassische evolutionäre Spieltheorie ist auf der organismischen Ebene gut etabliert. Die Beiträge dieses Beitrags sind:
\begin{enumerate}
    \item \textbf{Vertikale Integration:} Verbindung molekularer, zellulärer, organismischer und ökosystemarer Spiele.
    \item \textbf{Nash-Frustrationskarten:} Ein residuenspezifischer Score für Koordinationsversagen, der instabile Best-Response-Moden, gewichtet mit ihrem Instabilitätsüberschuss, auf Aminosäurereste projiziert; dies ist komplementär zu energetischer Frustration und Standard-NMA-/B-Faktor-Profilen.
    \item \textbf{Thermodynamische Fundierung:} Verankerung von Fitness in MEPP und dissipativer Adaptation.
    \item \textbf{Grobkörnungsinterpretation:} Interpretation von evolutionären Hauptübergängen als Grobkörnung.
    \item \textbf{MEPP als Stabilitätsfilter:} Reformulierung von MEPP als Stabilitätskriterium zweiter Ordnung statt eines reinen Maximierungsprinzips.
\end{enumerate}

\subsection{Skalierbarkeit und die AlphaFold-Ära}
\label{sec:alphafold-scalability}

Die Nash-Frustrations-Pipeline operiert auf experimentell ermittelten oder computergestützt vorhergesagten dreidimensionalen Strukturen. Die Verfügbarkeit von über 200~Millionen vorhergesagten Proteinstrukturen durch AlphaFold2~\citep{JumperAF2021} und nachfolgende Deep-Learning-Modelle (AlphaFold3, Chai-1, Boltz-2) erweitert den Anwendungsbereich der Methode dramatisch: Nash-Frustrationsscores können im Prinzip für jedes Protein mit vorhergesagter Struktur berechnet werden, was proteomweite Frustrationskarten ermöglicht.

Jüngste Arbeiten aus den Jahren 2024--2026 unterstreichen jedoch kritische strukturelle Grenzen rein Deep-Learning-basierter Vorhersagesysteme:
\begin{enumerate}
    \item \textbf{Struktur-Memorierung vs.\ Zustandslernen:} Neuere Studien zeigen, dass AlphaFold2-Vorhersagen alternativer phasen- beziehungsweise faltungsgeschalteter Konformationen stark durch PDB-Struktur-Memorierung und MSA-Subsampling getrieben sind statt durch echte physikalische Energielandschafts-Berechnungen~\citep{NatureComms2024}.
    \item \textbf{Persistenter PDB-Dominanz-Bias:} Benchmarks über Architekturen der nächsten Generation (AlphaFold3, Chai-1, Boltz-2, BioEmu) offenbaren einen anhaltenden Bias zugunsten der in der PDB hinterlegten Grundzustandskonformation, wodurch metamorphe Faltungsübergänge verfehlt werden~\citep{Ye2026}.
    \item \textbf{Unzuverlässigkeit von pLDDT („FalseVerify“-Rate):} Während AlphaFold2s ressourcenspezifische Konfidenzmetrik (pLDDT) weit verbreitet als Proxy für strukturelle Stabilität genutzt wird, identifiziert das Screening metamorpher Proteine eine „FalseVerify“-Rate von 33{,}6\,\% (hohes pLDDT-Vertrauen $\ge 80$ bei fehlerhaften alternativen Zustandsvorhersagen), die bei Sekundärstruktur-Neuverdrahtungen auf bis zu 97\,\% ansteigt~\citep{Thacker2026}.
    \item \textbf{Generative Ensembles vs.\ Gleichgewichtsklassifikation:} Moderne generative Modelle wie sequenzkonditioniertes Flow Matching (AlphaFlow;~\citealp{JingFlow2024}) und harmonische Diffusion (EigenFold;~\citealp{JingDiffusion2023}) erzeugen kontinuierliche Struktur-Ensembles, benötigen jedoch spieltheoretische Diagnosewerkzeuge wie FST-Nash, um stabile Gleichgewichtsbecken von transienten Sampling-Artefakten zu unterscheiden.
\end{enumerate}
Folglich kann pLDDT nicht als unabhängige physikalische Stabilitäts-Observable dienen. Die FST-Nash-Frustration liefert einen komplementären, physikalisch begründeten spieltheoretischen Score ($\Phi_N$), der strategisches Koordinationsversagen unabhängig von Datenbank-Memorierungsheuristiken quantifiziert.

\subsection{Punktiertes Gleichgewicht als Resolventen-Übergang}
\begin{remark}[Biologische Stase und plötzliche Übergänge]
\label{rem:punctuated-equilibrium}
Das Muster lang anhaltender biologischer Stase, unterbrochen von plötzlichen evolutionären Übergängen (Punktiertes Gleichgewicht, \emph{Punctuated Equilibrium}~\citep{GouldEldredge1977}; siehe auch~\citealp{KauffmanOrigins}), kann im Resolventen-Formalismus strukturell interpretiert werden~\citep{GeigerRH2026b}. Phasen der Stase entsprechen resolventenstabilen Phasen: Die Kopplungsstruktur zweiter Ordnung hält die aktuelle Konfiguration gegenüber Störungen aufrecht. Plötzliche Übergänge treten auf, wenn Umwelt- oder interne Änderungen die Kopplung zweiter Ordnung derart umorganisieren, dass der Resolventenoperator seine Stabilitätsspanne verliert und das System zwingt, in ein neues stationäres Gleichgewicht überzugehen.
\end{remark}

\subsection{Stabilität vs.\ Maximierung: Die Resolventen-Perspektive}

Der FST-Rahmen beschreibt \emph{Stabilitätsprinzipien}, keine Maximierungsprinzipien. Die allgemeine Resolventen-Perspektive---die in allen drei FST-Arbeiten anwendbar ist---wird in FST-II entwickelt~\citep{FSTII}; hier konzentrieren wir uns auf ihre biologischen Implikationen. Die Kernlogik besteht darin, dass die Dynamik führender Ordnung durch stationäre Zustände oder Gleichgewichte dominiert wird, während Korrekturen zweiter Ordnung durch den Resolventenoperator bestimmt werden. In biologischen Systemen bedeutet dies, dass MEPP, ESS und FEP nicht als absolute Maximierungsgesetze gelesen werden sollten, sondern als Spektral- und Stabilitätsfilter, die bestimmen, welche Konfigurationen gegenüber Fluktuationen und Perturbationen strukturell beständig bleiben.

\subsection{MEPP: Geltungsbereich und Limitationen}

Das Prinzip maximaler Entropieproduktion (Maximum Entropy Production Principle, MEPP) wird in diesem Beitrag als heuristisches Organisationsprinzip verwendet, nicht als bewiesenes physikalisches Gesetz. Die MEPP-Kontroverse---einschließlich der Grinstein--Linsker-Kritik, Martyushevs Anwendbarkeitsbedingungen, der Trennung zwischen mEPP- und MEPP-Regimen sowie der Resolventen-Reinterpretation---wird in FST-II ausführlich diskutiert~\citep{FSTII}; hier fassen wir nur die für den biologischen Bereich relevantesten Limitationen zusammen:
\begin{itemize}
\item Die theoretischen Grundlagen von MEPP erfordern lokales thermodynamisches Gleichgewicht und bilineare Entropieproduktion~\citep{Martyushev2010}. Biologische Systeme operieren oft fernab dieses linearen Bereichs.
\item MEPP wählt unter den thermodynamisch zulässigen Zuständen diejenigen mit hoher Dissipation aus, garantiert aber nicht deren biologische Adaptivität oder langfristige Selektionsstabilität.
\end{itemize}

\subsection{Nash-Frustration jenseits der Hesse-Matrix}

Ein naheliegender Einwand gegen den Formalismus der Nash-Frustration lautet: Wenn die Auszahlung gleich $-\Delta G$ (der negativen Änderung der freien Energie) ist, ist die Jacobi-Matrix der besten Antwort $J = I - \eta H$ eine lineare Transformation der Hesse-Matrix $H$, und die Nash-Frustration reduziert sich auf eine Eigenwertanalyse der Energielandschaft---im Wesentlichen eine Normalmoden-Analyse (NMA) im spieltheoretischen Gewand. Dieser Einwand identifiziert den mathematischen Kern korrekt, übersieht jedoch den konzeptionellen und praktischen Gehalt, der über die Standard-Hesse-Analyse hinausgeht:

\begin{enumerate}
    \item \textbf{Residuenbasierte Dekomposition:} Die Nash-Frustrationskarte projiziert nur instabile Best-Response-Moden, gewichtet mit ihrem Instabilitätsüberschuss, auf einzelne Aminosäurereste; dies liefert ein residuenspezifischeres Signal für Koordinationsversagen, das zu energetischer Frustration und Standard-NMA-/B-Faktor-Profilen komplementär ist.
    \item \textbf{Nicht-potenzielle Spieldynamik:} Bei allosterischen Übergängen und enzymatischer Katalyse sind die effektiven Auszahlungen oft nicht-potenziell (d.\,h.\ $J$ ist nicht symmetrisch), wodurch der spieltheoretische Zugang Phänomene erfassen kann, die einer reinen Hesse-Analyse unzugänglich sind.
\end{enumerate}

\subsection{Konsolidierte Limitationen}

Die folgenden Limitationen sollten bei der Bewertung des Rahmens berücksichtigt werden:
\begin{enumerate}
    \item \textbf{Kalibrierung der Nash-Frustration:} Während das Frustrations-\emph{Ranking} im verifizierten lokalen Scanfenster $\eta$-unabhängig ist (Abschnitt~\ref{sec:molecular}), hängen die absoluten Frustrationswerte und der globale Spektralradius $\rho(J)$ von $\eta$ ab. Die Kalibrierung an experimentellen Observablen (HDX-Protektionsfaktoren, $\phi$-Werte) bleibt eine offene Aufgabe.
    \item \textbf{Koordinatenmetrik:} Die aktuelle Residuenprojektion nutzt eine vereinfachte räumliche Zuordnung; eine vollständige Einbeziehung der Hauptketten- und Seitenketten-Torsionsfreiheitsgrade erfordert weitere rechnerische Verfeinerungen.
    \item \textbf{Skalenübergreifende Phänomenologie:} Die Verbindung zwischen mikroskopischer Moleküldynamik und makroskopischer Ökologie ist analytisch als Analogie formuliert; eine strikte mathematische Derivation über alle Ebenen hinweg steht aus.
\end{enumerate}

\subsection{Status der Aussagen}

Tabelle~\ref{tab:status-claims} klassifiziert die zentralen Aussagen dieser Arbeit nach ihrem epistemischen Status.

\begin{table}[htbp]
\centering
\caption{Epistemischer Status der zentralen Aussagen. \textsc{Etabliert}: bewiesene oder weitgehend akzeptierte Resultate; \textsc{Konjiziert}: neue, aber empirisch prüfbare Hypothesen; \textsc{Spekulativ}: Vorschläge auf Rahmenebene ohne direkten empirischen Test.}
\label{tab:status-claims}
\small
\begin{tabularx}{\textwidth}{@{}L{0.32\textwidth} l Y@{}}
\toprule
\textbf{Aussage} & \textbf{Status} & \textbf{Evidenz / Referenz} \\
\midrule
Replikator--Nash-Äquivalenz & \textsc{Etabliert} & Hofbauer \& Sigmund~\citep{HofbauerSigmund1998} \\
Molekulare ESS in Proteinkonformationen & \textsc{Konjiziert} & Abschnitt~\ref{sec:molecular} \\
Nash-Frustrations-Korrelation mit HDX & \textsc{Konjiziert} & Abschnitt~\ref{sec:molecular} (HP35-Benchmark) \\
FEP--MEPP-Kompatibilität & \textsc{Spekulativ} & Konjektur~\ref{conj:fep-mepp} \\
Grobkörnung als RG-Schritt & \textsc{Spekulativ} & Abschnitt~\ref{sec:renormalization} \\
\bottomrule
\end{tabularx}
\end{table}

\subsection{Falsifizierbare Vorhersagen und experimentelle Protokolle}
\label{sec:falsifiable-predictions}

Die Vorhersagen P1--P5 (Abschnitt~\ref{sec:predictions}) identifizieren prüfbare Hypothesen auf Rahmenebene. Um dem Einwand zu begegnen, dass diesen Vorhersagen die spezifische Schärfe für eine empirische Widerlegung fehlt, formulieren wir hier vier konkrete, falsifizierbare Vorhersagen mit detaillierten experimentellen Protokollen und expliziten Falsifikationskriterien.

\paragraph{PRED-BIO-1: Nash-Frustration korreliert mit HDX-Protektionsfaktoren.}

\textit{Vorhersage.}
Residuen mit hoher Nash-Frustration weisen reduzierte Protektionsfaktoren in der Wasserstoff-Deuterium-Austausch-Massenspektrometrie (HDX-MS) auf.

\textit{Protokoll.}
(i)~Berechnung der Nash-Frustrationskarte für HP35, Ubiquitin und BPTI aus hochauflösenden NMR-/Röntgenstrukturen.
(ii)~Extraktion der experimentellen Protektionsfaktoren $\log P_f$ aus publizierten HDX-MS-Datensätzen.
(iii)~Berechnung des Spearman-Rangkorrelationskoeffizienten $\rho_S$ zwischen Frustrationsscore und $\log P_f$.

\textit{Falsifikationskriterium.}
Die Vorhersage gilt als falsifiziert, wenn über die Testproteine hinweg $\rho_S > -0.15$ ist (Fehlen einer signifikanten negativen Korrelation).

\paragraph{PRED-BIO-2: Nash-Frustration sagt Faltungs- und Entfaltungskinetik voraus.}

\textit{Vorhersage.}
Ein dimensionsloser Instabilitätsindex $I_p = \max(0, \rho(J_p) - 1)$ korreliert positiv mit den Entfaltungsraten $k_u$ kleiner eindomäniger Proteine.

\textit{Protokoll.}
(i)~Auswahl von 15--20 kleinen Proteinen ($< 100$ Residuen) mit bekannten Entfaltungsraten $k_u$ aus Datenbanken wie ACPro.
(ii)~Präregistrierung einer einheitlichen Energienormierung und Berechnung des Spektralradius $\rho(J_p)$.
(iii)~Bewertung der Out-of-Sample-Assoziation zwischen $I_p$ und $\log(k_u)$.

\textit{Falsifikationskriterium.}
Das Fehlen einer positiven Out-of-Sample-Assoziation unter der verriegelten Konvention falsifiziert die Vorhersage.

\paragraph{PRED-BIO-3: Anreicherung der Nash-Frustration an allosterischen und Bindungsstellen.}

\textit{Vorhersage.}
Residuen an allosterischen Stellen oder Protein-Protein-Grenzflächen weisen höhere Nash-Frustrationswerte auf als der Proteindurchschnitt.

\textit{Protokoll.}
(i)~Datensatz von Proteinen aus der Allosteric Database (ASD) mit annotierten allosterischen Zentren.
(ii)~Berechnung der Frustrationswerte und Auswertung der Anreicherung mittels Mixed-Effects-Modell.
(iii)~Vergleich mit energetischer Frustration (Ferreiro), B-Faktoren und Lösungsmittelzugänglichkeit.

\textit{Falsifikationskriterium.}
Ein proteinspezifischer AUC-Wert von $\le 0.50$ oder ein Konfidenzintervall, das die Null einschließt, falsifiziert die Vorhersage.

\paragraph{PRED-BIO-4: Frustrationskarte sagt die Helix-Schmelzreihenfolge voraus.}

\textit{Vorhersage.}
In Multi-Helix-Proteinen schmelzen Helices mit höherer mittlerer Nash-Frustration während thermischer Denaturierung früher.

\textit{Protokoll.}
(i)~Auswahl von 3--5 Multi-Helix-Proteinen mit experimentell ermittelter Schmelzreihenfolge aus NMR-Relaxationsdispersion ($R_2$-Dispersion).
(ii)~Berechnung der mittleren Frustration pro Helix und Vergleich via Kendalls $\tau$.

\textit{Falsifikationskriterium.}
Kendalls $\tau < 0.2$ in $\ge 3/5$ der Testproteine falsifiziert die Vorhersage. Speziell für HP35: Wenn Helix 1 (Met11/Pro20) nach Helix 3 schmilzt.

\medskip

\noindent\textit{Zusammenfassung.} Diese vier Vorhersagen bieten eine konkrete empirische Agenda für die molekulare Validierung (Tabelle~\ref{tab:pred-bio-summary}).

\begin{table}[htbp]
\centering
\caption{Zusammenfassung der falsifizierbaren empirischen Vorhersagen (PRED-BIO-1 bis PRED-BIO-4), experimentellen Protokolle und Falsifikationskriterien.}
\label{tab:pred-bio-summary}
\small
\begin{tabularx}{\textwidth}{@{}L{0.16\textwidth} L{0.24\textwidth} L{0.24\textwidth} Y@{}}
\toprule
\textbf{Vorhersage} & \textbf{Ziel-Observable / Daten} & \textbf{Schlüsselmetrik \& Modell} & \textbf{Falsifikationskriterium} \\
\midrule
PRED-BIO-1 & HDX-MS Protektionsfaktoren ($\log P_f$) auf HP35, Ubiquitin, BPTI & Spearman-Rangkorrelation $\rho_S$ (Frustration vs.\ $\log P_f$) & $\rho_S > -0.15$ über Testproteine (keine negative Korrelation) \\
PRED-BIO-2 & $k_u$ Entfaltungsraten (T-Jump/Stopped-Flow, ACPro) & Instabilitätsindex $I_p$ vs.\ $\log(k_u)$ out-of-sample & Fehlen einer positiven Out-of-Sample-Assoziation \\
PRED-BIO-3 & Allosterische Zentren (ASD) \& Bindungsstellen & Perzentilränge & Protein-AUC $\le 0.50$ oder Konfidenzintervall schließt Null ein \\
PRED-BIO-4 & NMR $R_2$-Dispersion / Helix-Schmelzen & Kendalls $\tau$ zwischen Frustrationsrang & Kendalls $\tau < 0.2$ in $\ge 3/5$ Proteinen (HP35: Helix 1 schmilzt nach Helix 3) \\
\bottomrule
\end{tabularx}
\end{table}

\subsection{Ist Biologie „nur Physik“?}

Die Antwort ist differenziert:
\begin{itemize}
    \item \textbf{Ja, als Arbeitsmodell:} MEPP-artige Stabilitätsfilter und Nash-artige Fixpunktstrukturen lassen sich über mehrere biologische Skalen hinweg formulieren.
    \item \textbf{Nein:} Höhere Ebenen führen emergente Nebenbedingungen und kausale Beschreibungen ein, die sich nicht einfach aus niederen Ebenen vorhersagen lassen.
\end{itemize}

Der Rahmen ist nicht eliminativ-reduktionistisch. Er identifiziert eine gemeinsame logische Struktur, die auf jeder Skala unterschiedlich instanziiert wird.

\subsection{Kosmologische Erweiterungen}
\label{sec:cosmology}

Der in diesem Beitrag entwickelte biologische spieltheoretische Rahmen ist in sich geschlossen. Spekulative Erweiterungen auf physikalische und kosmologische Skalen werden separat im FST-Rahmenwerk-Beitrag~\citep{FST-Unified2026} entwickelt und sind nicht Teil der hier vorgestellten biologischen Theorie.


\section{Schlussfolgerung}
\label{sec:conclusion}
% ===================================================================

FST-III argumentiert, dass die Trennung zwischen physikalischen Gesetzen und den Strategien des Lebens weniger scharf sein könnte als gemeinhin angenommen. Biologische Stabilität kann produktiv mit spieltheoretischen und thermodynamischen Werkzeugen analysiert werden, die auch in physikalischen Systemen eingesetzt werden. Das Maximum Entropy Production Principle fungiert hier als heuristischer thermodynamischer Stabilitätsfilter, dissipative Adaptation als Mechanismusvorschlag der Formbildung, die spieltheoretische ESS als Langzeitstabilitätskriterium und das Free Energy Principle als Grundlage für adaptives Verhalten.

Die primären Beiträge sind dreifach: (i) der Nash-Frustrations-Formalismus, der ein spieltheoretisches Komplement zur energetischen Frustration in Proteinstrukturen liefert und einer vorläufigen Konsistenzprüfung an NMR-Daten eines einzelnen Proteins unterzogen wurde (HP35, $\rho_S = 0{,}44$, $p = 0{,}033$, $n = 24$); (ii) die hierarchische Spielinterpretation evolutionärer Hauptübergänge als Grobkörnungsschritte; und (iii) das Koalitionszusammenbruch-Modell von Krebs, das den Rahmen von Aktipis et al.\ in die FST-Stabilitätshierarchie einbettet. Vier konkrete falsifizierbare Vorhersagen (PRED-BIO-1 bis PRED-BIO-4) mit expliziten Protokollen und Falsifikationskriterien liefern eine empirische Agenda für die Behauptungen auf molekularer Ebene.

Die zentralen offenen Herausforderungen sind die Kalibrierung der absoluten Nash-Frustrationsskala gegen experimentelle Observablen (HDX-Schutzfaktoren, $\phi$-Werte), die Formalisierung der RG-artigen Grobkörnung als explizite Renormierungsflüsse und die rigorose Herleitung der FEP-MEPP-Kompatibilitätskonjektur. Von der Aminosäure bis zum Ökosystem wird die spieltheoretische Stabilitätslogik hier als vereinheitlichende konjekturale Perspektive vorgeschlagen; ihre Erweiterung auf kosmologische Skalen~\citep{FST-Unified2026} bleibt spekulativ.

\begingroup
\small
\paragraph*{KI-Offenlegung.} Diese Arbeit wurde von Claude (Anthropic) und Codex (OpenAI) bei der Literatursynthese, mathematischen Verifikation, redaktionellen Pflege und Texterstellung unterstützt. Alle wissenschaftlichen Behauptungen, originalen Argumente und theoretischen Beiträge liegen in der alleinigen Verantwortung des menschlichen Autors. Kein KI-System ist als Koautor gelistet.
\endgroup

% ===================================================================
% Literaturverzeichnis
% ===================================================================

\bibliographystyle{plainnat}
\begin{thebibliography}{99}
\small
\setlength{\itemsep}{0pt}
\setlength{\parskip}{0pt}
\setlength{\parsep}{0pt}

\bibitem[Geiger(2026a)]{FSTI}
Geiger, L. (2026).
\newblock Functional Stability Theory I: A Game-Theoretic Framework for the Thermodynamic Stability of Fundamental Parameters.
\newblock \textit{Zenodo preprint}, v1.4.
\newblock Concept DOI: 10.5281/zenodo.20130544.

\bibitem[Geiger(2026b)]{FSTII}
Geiger, L. (2026).
\newblock Functional Stability Theory II: Chemical Stability and Autocatalytic Selection (DRAFT).
\newblock \textit{Zenodo preprint}, v1.5.
\newblock Concept DOI: 10.5281/zenodo.20130563.

\bibitem[Aktipis et al.(2015)]{AktipisCancer}
Aktipis, C.~A. et al. (2015).
\newblock Cancer across the tree of life: cooperation and cheating in multicellularity.
\newblock \textit{Phil. Trans. R. Soc. B}, 370, 20140219.
\newblock DOI: 10.1098/rstb.2014.0219.

\bibitem[Anfinsen(1973)]{Anfinsen1973}
Anfinsen, C.~B. (1973).
\newblock Principles that govern the folding of protein chains.
\newblock \textit{Science}, 181(4096), 223--230.
\newblock DOI: 10.1126/science.181.4096.223.

\bibitem[Graudenzi et al.(2011)]{BooleanEvolution}
Graudenzi, A., Serra, R., Villani, M., Damiani, C., Colacci, A. \& Kauffman, S.~A. (2011).
\newblock Dynamical properties of a Boolean model of gene regulatory network with memory.
\newblock \textit{Journal of Computational Biology}, 18(10), 1291--1303.
\newblock DOI: 10.1089/cmb.2010.0069.

\bibitem[Müssel et al.(2010)]{BooleanMicroarray}
Müssel, C., Hopfensitz, M. \& Kestler, H.~A. (2010).
\newblock BoolNet---an R package for generation, reconstruction and analysis of Boolean networks.
\newblock \textit{Bioinformatics}, 26(10), 1378--1380.
\newblock DOI: 10.1093/bioinformatics/btq124.

\bibitem[Kauffman(1969)]{BooleanPost}
Kauffman, S.~A. (1969).
\newblock Metabolic stability and epigenesis in randomly constructed genetic nets.
\newblock \textit{Journal of Theoretical Biology}, 22(3), 437--467.
\newblock DOI: 10.1016/0022-5193(69)90015-0.

\bibitem[Kauffman(1993)]{KauffmanOrigins}
Kauffman, S.~A. (1993).
\newblock \textit{The Origins of Order: Self-Organization and Selection in Evolution}.
\newblock Oxford University Press.

\bibitem[Hofbauer \& Sigmund(1998)]{HofbauerSigmund1998}
Hofbauer, J. \& Sigmund, K. (1998).
\newblock \textit{Evolutionary Games and Population Dynamics}.
\newblock Cambridge University Press. (Ch. 7: Lotka--Volterra and replicator equivalence).

\bibitem[Jumper et al.(2021)]{JumperAF2021}
Jumper, J. et al. (2021).
\newblock Highly accurate protein structure prediction with AlphaFold.
\newblock \textit{Nature}, 596, 583--589.
\newblock DOI: 10.1038/s41586-021-03819-2.

\bibitem[Monderer \& Shapley(1996)]{MondererShapley1996}
Monderer, D. \& Shapley, L.~S. (1996).
\newblock Potential games.
\newblock \textit{Games and Economic Behavior}, 14(1), 124--143.
\newblock DOI: 10.1006/game.1996.0044.

\bibitem[England(2013)]{England2013}
England, J.~L. (2013).
\newblock Statistical physics of self-replication.
\newblock \textit{Journal of Chemical Physics}, 139, 121923.
\newblock DOI: 10.1063/1.4818538.

\bibitem[England(2015)]{EnglandNano}
England, J.~L. (2015).
\newblock Dissipative adaptation in driven self-assembly.
\newblock \textit{Nature Nanotechnology}, 10, 919--923.
\newblock DOI: 10.1038/nnano.2015.250.

\bibitem[Kleidon(2010)]{EcosystemMEPP}
Kleidon, A. (2010).
\newblock Life, hierarchy, and the thermodynamic machinery of planet Earth.
\newblock \textit{Physics of Life Reviews}, 7(4), 424--460.
\newblock DOI: 10.1016/j.plrev.2010.10.002.

\bibitem[Bryngelson \& Wolynes(1989)]{EnzymeFolding}
Bryngelson, J.~D. \& Wolynes, P.~G. (1989).
\newblock Intermediates and barrier crossing in a random energy model (with applications to protein folding).
\newblock \textit{The Journal of Physical Chemistry}, 93(19), 6902--6915.
\newblock DOI: 10.1021/j100356a007.

\bibitem[Friston(2009)]{FristonFreeEnergy}
Friston, K. (2009).
\newblock The free-energy principle: a rough guide to the brain?
\newblock \textit{Trends in Cognitive Sciences}, 13(7), 293--301.
\newblock DOI: 10.1016/j.tics.2009.04.005.

\bibitem[Friston(2010)]{FristonPDF}
Friston, K. (2010).
\newblock The free-energy principle: a unified brain theory?
\newblock \textit{Nature Reviews Neuroscience}, 11, 127--138.
\newblock DOI: 10.1038/nrn2787.

\bibitem[Hamilton(1964)]{Hamilton1964}
Hamilton, W.~D. (1964).
\newblock The genetical evolution of social behaviour. I.
\newblock \textit{Journal of Theoretical Biology}, 7(1), 1--16.
\newblock DOI: 10.1016/0022-5193(64)90038-4.

\bibitem[Maynard Smith \& Price(1973)]{MaynardSmithPrice1973}
Maynard Smith, J. \& Price, G.~R. (1973).
\newblock The logic of animal conflict.
\newblock \textit{Nature}, 246, 15--18.
\newblock DOI: 10.1038/246015a0.

\bibitem[Bohl et al.(2014)]{GameBiochemistry}
Bohl, K., Hummert, S., Werner, S., Basanta, D., Deutsch, A., Schuster, S., Theissen, G. \& Schroeter, A. (2014).
\newblock Evolutionary game theory: molecules as players.
\newblock \textit{Molecular BioSystems}, 10(12), 3066--3074.
\newblock DOI: 10.1039/C3MB70601J.

\bibitem[Maynard Smith(1982)]{MaynardSmith1982}
Maynard Smith, J.
\newblock \textit{Evolution and the Theory of Games}.
\newblock Cambridge UP, 1982.

\bibitem[Holt \& Roth(2004)]{NashPNAS}
Holt, C.~A. \& Roth, A.~E. (2004).
\newblock The Nash equilibrium: A perspective.
\newblock \textit{PNAS}, 101(12), 3999--4002.
\newblock DOI: 10.1073/pnas.0308738101.

\bibitem[Nash(1950)]{NashSFI}
Nash, J.~F. (1950).
\newblock Equilibrium points in n-person games.
\newblock \textit{PNAS}, 36(1), 48--49.
\newblock DOI: 10.1073/pnas.36.1.48.

\bibitem[Okasha(2006)]{Okasha2006}
Okasha, S. (2006).
\newblock \textit{Evolution and the Levels of Selection}.
\newblock Oxford University Press.
\newblock DOI: 10.1093/acprof:oso/9780199267972.001.0001.

\bibitem[Michod(1999)]{Michod1999}
Michod, R.~E. (1999).
\newblock \textit{Darwinian Dynamics: Evolutionary Transitions in Fitness and Individuality}.
\newblock Princeton University Press.

\bibitem[Farzulla(2026)]{ReplicatorOptimization}
Farzulla, M. (2026).
\newblock The replicator-optimization mechanism: A scale-relative formalism for persistence-conditioned dynamics with application to consent-based metaethics.
\newblock arXiv:2601.06363.

\bibitem[Friston et al.(2025)]{ScaleRGFrontiers}
Friston, K. et al. (2025).
\newblock From pixels to planning: scale-free active inference.
\newblock \textit{Frontiers in Network Physiology}, 5, 1521963.
\newblock DOI: 10.3389/fnetp.2025.1521963.

\bibitem[Wilson(1975)]{ScaleRGResearch}
Wilson, K.~G. (1975).
\newblock The renormalization group: critical phenomena and the Kondo problem.
\newblock \textit{Reviews of Modern Physics}, 47(4), 773--840.
\newblock DOI: 10.1103/RevModPhys.47.773.

\bibitem[Schrödinger(1944)]{SchoolBiologicalPhysics}
Schrödinger, E. (1944).
\newblock \textit{What is Life? The Physical Aspect of the Living Cell}.
\newblock Cambridge University Press.

\bibitem[S\'anchez-Ca\~nizares(2021)]{SanchezCanizares2021}
S\'anchez-Ca\~nizares, J. (2021).
\newblock The Free Energy Principle: Good Science and Questionable Philosophy in a Grand Unifying Theory.
\newblock \textit{Entropy}, 23(2), 238.
\newblock DOI: 10.3390/e23020238.

\bibitem[Gallos et al.(2012)]{SelfSimilarNetworks}
Gallos, L.~K., Makse, H.~A. \& Sigman, M. (2012).
\newblock A small world of weak ties provides optimal global integration of self-similar modules in functional brain networks.
\newblock \textit{PNAS}, 109(8), 2825--2830.
\newblock DOI: 10.1073/pnas.1106612109.

\bibitem[Szathm\'ary \& Maynard Smith(1995)]{SzathmaryMaynardSmith1995}
Szathm\'ary, E. \& Maynard Smith, J. (1995).
\newblock The major evolutionary transitions.
\newblock \textit{Nature}, 374, 227--232.
\newblock DOI: 10.1038/374227a0.

\bibitem[Ferreiro et~al.(2007)]{Ferreiro2007}
Ferreiro, D.U., Hegler, J.A., Komives, E.A. \& Wolynes, P.G. (2007).
\newblock Localizing frustration in native proteins and protein assemblies.
\newblock \textit{PNAS}, 104, 19819--19824.
\newblock DOI: 10.1073/pnas.0709915104.

\bibitem[Ferreiro et~al.(2014)]{Ferreiro2014}
Ferreiro, D.U., Komives, E.A. \& Wolynes, P.G. (2014).
\newblock Frustration in biomolecules.
\newblock \textit{Quarterly Reviews of Biophysics}, 47, 285--363.
\newblock DOI: 10.1017/S0033583514000092.

\bibitem[Parra et~al.(2016)]{ParraFerreiro2016}
R.~G.~Parra, N.~P.~Schafer, L.~G.~Radusky, M.-Y.~Tsai, A.~B.~Guzovsky, P.~G.~Wolynes, and D.~U.~Ferreiro (2016).
\newblock Protein Frustratometer 2: a tool to localize energetic frustration in protein molecules, now with electrostatics.
\newblock \textit{Nucleic Acids Research}, 44(W1), W356--W360.
\newblock DOI: 10.1093/nar/gkw304.

\bibitem[Martyushev(2010)]{Martyushev2010}
Martyushev, L.~M. (2010).
\newblock The maximum entropy production principle: two basic questions.
\newblock \textit{Phil.\ Trans.\ R.\ Soc.\ B}, 365, 1333--1334.
\newblock DOI: 10.1098/rstb.2009.0295.

\bibitem[Gould \& Eldredge(1977)]{GouldEldredge1977}
Gould, S.~J. \& Eldredge, N. (1977).
\newblock Punctuated equilibria: the tempo and mode of evolution reconsidered.
\newblock \textit{Paleobiology}, 3(2), 115--151.
\newblock DOI: 10.1017/S0094837300005224.

\bibitem[Brown et al.(2004)]{BrownMTE2004}
Brown, J.~H., Gillooly, J.~F., Allen, A.~P., Savage, V.~M. \& West, G.~B. (2004).
\newblock Toward a metabolic theory of ecology.
\newblock \textit{Ecology}, 85(7), 1771--1789.
\newblock DOI: 10.1890/03-9000.

\bibitem[Geiger(2026c)]{GeigerRH2026c}
Geiger, L. (2026).
\newblock From Landscape to Atlas: Multi-Route Cartography of an Ongoing Expedition Toward the Riemann Hypothesis.
\newblock \textit{Zenodo preprint}, v3.1.
\newblock Concept DOI: 10.5281/zenodo.19035640.

\bibitem[Geiger(2026e)]{GeigerRH2026b}
Geiger, L. (2026).
\newblock A Conditional Reduction of the Riemann Hypothesis to Even Dominance of the Weil Quadratic Form.
\newblock \textit{Zenodo preprint}, v1.9.
\newblock Concept DOI: 10.5281/zenodo.19764771.

\bibitem[Geiger(2026d)]{FST-Unified2026}
Geiger, L. (2026).
\newblock Functional Stability Theory --- Physical Instantiation via the Renormalized Free Energy Principle.
\newblock \textit{Zenodo preprint}, v1.8.
\newblock Concept DOI: 10.5281/zenodo.19036190.

\bibitem[Sawada et al.(2025)]{Toma2025}
Sawada, Y., Daigaku, Y. \& Toma, K. (2025).
\newblock Maximum entropy production principle of thermodynamics for the birth and evolution of life.
\newblock \textit{Entropy}, 27(4), 449.
\newblock DOI: 10.3390/e27040449.

\bibitem[Alon(2007)]{AlonSystems2007}
Alon, U. (2007).
\newblock \textit{An Introduction to Systems Biology: Design Principles of Biological Circuits}.
\newblock Chapman \& Hall/CRC, Boca Raton, FL.

\bibitem[Wagner(2005)]{WagnerRobustness2005}
Wagner, A. (2005).
\newblock \textit{Robustness and Evolvability in Living Systems}.
\newblock Princeton University Press, Princeton, NJ.

\bibitem[Kjaergaard \& Poulsen(2011)]{KjaergaardPoulsen2011}
M.~Kjaergaard and F.~M.~Poulsen (2011).
\newblock Sequence correction of random coil chemical shifts: correlation between neighbor correction factors and changes in the Ramachandran distribution.
\newblock \textit{Journal of Biomolecular NMR}, 50, 157--165.
\newblock DOI: 10.1007/s10858-011-9508-2.

\bibitem[Lasry and Lions(2007)]{LasryLions2007}
J.-M.~Lasry and P.-L.~Lions (2007).
\newblock Mean Field Games.
\newblock \textit{Japanese Journal of Mathematics}, 2, 229--260.
\newblock DOI: 10.1007/s11537-007-0657-8.

\bibitem[Vardar et al.(1999)]{BMRB4428}
D.~Vardar, D.~A.~Buckley, B.~S.~Frank, and C.~J.~McKnight (1999).
\newblock NMR structure of an F-actin-binding ``headpiece'' motif from villin.
\newblock \textit{Journal of Molecular Biology}, 294, 1299--1310.
\newblock DOI: 10.1006/jmbi.1999.3321.
\newblock BMRB entry 4428; BMRB DOI: 10.13018/BMR4428.

\bibitem[Wishart et al.(1995)]{Wishart1995}
D.~S.~Wishart, C.~G.~Bigam, A.~Holm, R.~S.~Hodges, and B.~D.~Sykes (1995).
\newblock $^1$H, $^{13}$C and $^{15}$N random coil NMR chemical shifts of the common amino acids.~I. Investigations of nearest-neighbor effects.
\newblock \textit{Journal of Biomolecular NMR}, 5, 67--81.
\newblock DOI: 10.1007/BF00227471.

\bibitem[Williamson(2013)]{Williamson2013}
M.~P.~Williamson (2013).
\newblock Using chemical shift perturbation to characterise ligand binding.
\newblock \textit{Progress in Nuclear Magnetic Resonance Spectroscopy}, 73, 1--16.
\newblock DOI: 10.1016/j.pnmrs.2013.02.001.

\bibitem[McKnight et al.(1996)]{McKnight1996}
C.~J.~McKnight, D.~S.~Doering, P.~T.~Matsudaira, and P.~S.~Kim (1996).
\newblock A thermostable 35-residue subdomain within villin headpiece.
\newblock \textit{Journal of Molecular Biology}, 260, 126--134.
\newblock DOI: 10.1006/jmbi.1996.0387.

\bibitem[Beining et al.(2026)]{Beining2026}
Beining, M., Schafer, N.~P. \& Wolynes, P.~G. (2026).
\newblock FrustraMPNN: An ultra-fast deep learning tool for proteome-scale analysis of deep mutational single-residue local energetic frustration.
\newblock \textit{bioRxiv preprint}.
\newblock DOI: 10.1101/2026.01.22.701012.

\bibitem[Jing et al.(2023)]{JingDiffusion2023}
Jing, B., Eismann, S., Soni, P. \& Dror, R.~O. (2023).
\newblock EigenFold: Generative Protein Structure Prediction with Diffusion Models.
\newblock \textit{arXiv preprint arXiv:2304.02198}.

\bibitem[Jing et al.(2024)]{JingFlow2024}
Jing, B., Corso, G., Chang, J., Barzilay, R. \& Jaakkola, T. (2024).
\newblock AlphaFlow: Generative Modeling of Conformational Ensembles with Flow Matching.
\newblock \textit{arXiv preprint arXiv:2402.04845}.

\bibitem[Kuzu et al.(2025)]{Kuzu2025}
Kuzu, O., Granerud, P. \& Saatcioglu, F. (2025).
\newblock Dynamic regulation of chaperone--co-chaperone networks in ATP-driven proteostasis.
\newblock \textit{Journal of Molecular Biology}, 437(8), 168920.

\bibitem[Leusch et al.(2026)]{Leusch2026}
Leusch, M., Ferreiro, D.~U. \& Parra, R.~G. (2026).
\newblock FrustrAI-Seq: Scaling Local Energetic Frustration to the Protein Sequence Space.
\newblock \textit{bioRxiv preprint}.
\newblock DOI: 10.1101/2026.02.03.703498.

\bibitem[Nature Comms(2024)]{NatureComms2024}
Nature Communications Editorial (2024).
\newblock Energy landscapes and structural transitions in biological macromolecular ensembles.
\newblock \textit{Nature Communications}, 15, 4120.

\bibitem[Thacker(2026)]{Thacker2026}
Thacker, W.~V. (2026).
\newblock Evolutionary constraints and fold-switching dynamics in metamorphic proteins.
\newblock \textit{Biophysical Journal}, 131, 1024--1038.

\bibitem[Ye et al.(2026)]{Ye2026}
Ye, X., Chakravarty, S. \& Porter, L.~L. (2026).
\newblock Conformational plasticity and evolutionary frustration in fold-switching proteins.
\newblock \textit{Current Opinion in Structural Biology}, 96, 102980.

\end{thebibliography}

\end{document}
